\section{Levels, boundary ends, and projective comparison}
\label{sec:D}

We use $\Gamma_{g,n}$ for the orientation-preserving mapping class group fixing every labelled puncture, $d=3g-3+n$, and $\mathcal C_{g,n}$ for the unmodified disjointness curve complex. For $0\leq k<d$, $P_k$ is its face poset restricted to multicurves of cardinality at least $d-k$. For $k\geq d$ we include the empty multicurve. This convention is essential at complexity one. The geometric realization of the empty-multicurve extension is a cone.

\subsection{The actual end and its fundamental group}

\begin{result}{D.1}{Proved}{The end comparison, including the exceptional types}
Let $2g-2+n>0$. Let $L_{g,n}$ be the boundary of the real oriented blow-up of the analytic Deligne--Mumford stack $\overline{\mathcal M}_{g,n}$, with corners rounded. The natural map
\[
 L_{g,n}\longrightarrow\mathcal M_{g,n}^{\rm an}\simeq B\Gamma_{g,n}
\]
has homotopy fibre $|\mathcal C_{g,n}|$. For nonempty boundary it induces an equivalence of fundamental groupoids precisely in the following ranges:
\[
 g=0,\ n\geq6;\qquad g=1,\ n\geq3;\qquad g\geq2,\ n\geq0.
\]
For $(0,5)$ and $(1,2)$ it induces a surjection on fundamental groups with kernel a free group of countably infinite rank. For $(0,4)$ and $(1,1)$ it does not induce a surjection on the fundamental group at a fixed boundary component. At $(0,3)$ the boundary is empty.
\end{result}
\begin{proof}
Choose an invariant thin part of Teichm\"uller space: a positive length threshold sufficiently small that curves shorter than the threshold are pairwise disjoint. The collar lemma gives this choice. For a nonempty multicurve $\sigma$, the locus where all its curves are short is contractible by Fenchel--Nielsen length and twist coordinates; intersections are the corresponding loci for unions of compatible multicurves, and are empty otherwise. The nerve of this cover is $\mathcal C_{g,n}$. A partition of unity provides the equivariant nerve comparison; equivariance can alternatively be obtained from the homotopy colimit of the cover, avoiding a choice of invariant partition on an infinite complex. The thin part is an equivariant model for the boundary collar of Harvey's bordification. Consequently the stack end is
\[
 E\Gamma_{g,n}\times_{\Gamma_{g,n}}|\mathcal C_{g,n}|,
\]
and its projection to $B\Gamma_{g,n}$ is the claimed homotopy-fibre sequence. This is the comparison used in EXIT\_ONE\_TYPE P5b1; it retains the stabilizer classifying spaces.

The exact sphere dimensions are
\[
 h(g,n)=\begin{cases}n-4&g=0,\\2g-3+n&g\geq1,\ n\geq1,\\2g-2&g\geq2,\ n=0.\end{cases}
\]
Harer's sphere theorem, in the explicit formulation of Birman--Broaddus--Menasco, Known Result~5 and equation~(1) \cite{BBM2015}, identifies the curve complex with a countably infinite wedge of $h(g,n)$-spheres. We have opened that formulation; Harer's original subscription text was not available. For $h\geq2$, the long exact sequence of the displayed fibre sequence gives connectedness and the asserted isomorphism of fundamental groups. For $h=1$, it gives
\[
 1\longrightarrow F_{\aleph_0}\longrightarrow\pi_1(L_{g,n})
 \longrightarrow\Gamma_{g,n}\longrightarrow1,
\]
since $\pi_2(B\Gamma_{g,n})=0$. For $h=0$, choose a curve $c$. Its component has group $\operatorname{Stab}_{\Gamma_{g,n}}(c)$, and the map is the stabilizer inclusion. The orbit of $c$ is infinite, so this subgroup is proper. These cases are exactly $(0,4)$ and $(1,1)$. At $(0,3)$ there are no stable graphs with an edge. Substituting the displayed dimensions proves the complete list.
\end{proof}

\begin{result}{D.2}{Proved}{Depth three is the sharp uniform input for the end's one-type}
In the normal-crossings link construction, retain every deployed stratum of depth at most $r$ and every flag among them. The resulting homotopy colimit maps naturally to the full link by an $(r-1)$-connected map. Thus depths $1,2,3$ compute its fundamental groupoid. Depths $1,2$ do not suffice uniformly: at the transverse meeting of three branches the fundamental groups are $\mathbb Z^4$ for the truncated diagram and $\mathbb Z^3$ for the full link. Applied to the stable moduli stack, this computes $\Pi_1(L_{g,n})$; D.1 gives precisely the types where it computes $\Pi_1(\mathcal M_{g,n})$. It does not compute $\Pi_1(\overline{\mathcal M}_{g,n})$ by that identification.
\end{result}
\begin{proof}
At a point on $q$ local branches, the normal angular factor is $(S^1)^q$. The auxiliary flag direction for the truncated diagram is the order complex of nonempty subsets of $\{1,\ldots,q\}$ of size at most $r$. This is the barycentric subdivision of the $(r-1)$-skeleton of a $(q-1)$-simplex. Its inclusion in the simplex is $(r-1)$-connected: the omitted cells have dimension at least $r$, and the skeleton is $(r-2)$-connected. Taking its product with $(S^1)^q$ gives the local comparison, including its homotopy fibre. The local contractions and their face maps are natural under permutations of the branches. They therefore descend through the finite stabilizer charts of the moduli stack and glue in the deployed flag homotopy colimit. Relative cellular connectivity is preserved under these gluings. For $q=3,r=2$ the auxiliary complex is the boundary of a triangle. Thus the local truncated model is $(S^1)^3\times S^1$, whereas the full auxiliary simplex is contractible. The displayed groups follow by the product formula. The map kills the last factor and is the identity on the three normal meridians. This proves both sharpness and the exact corrected target of cote~147, scan~152, formula~(247).
\end{proof}

\begin{result}{D.3}{Proved}{Compactification independence of the two end constructions}
For any compact Hausdorff compactification $U\subset X$ with $D=X\setminus U$, the archimedean punctured-neighborhood pro-space is canonically the end system $\{U\setminus K\}_{K\subset U\text{ compact}}$. For smooth proper normal-crossings pairs this pro-homotopy type is constant, with representative the rounded real-oriented-blow-up boundary. Over a trivially valued field $k$, a proper modification $X'\to X$ equal to the identity on $U$ induces an isomorphism
\[
 W(X',X'\setminus U)\xrightarrow{\ \sim\ }W(X,D),\qquad
 W(X,D)=X^{\beth}\setminus(U^{\beth}\cup D^{\beth}).
\]
For proper $X$, this space is $U^{\rm an}\setminus U^{\beth}$. Properness cannot be removed: the compactifications $\mathbb G_m\subset\mathbb A^1$ and $\mathbb G_m\subset\mathbb P^1$ give one and two punctured boundary tubes respectively.
\end{result}
\begin{proof}
The mutually inverse assignments are $N\mapsto K=X\setminus N$ and $K\mapsto N=X\setminus K$, where $N$ is an open neighborhood of $D$. Compact Hausdorffness makes each $K$ compact and each embedded compact $K$ closed. They identify the actual punctured open sets and every refinement map. In the normal-crossings chart $z_1\cdots z_q=0$, the oriented blow-up has coordinates $z_i=\rho_i e^{i\theta_i}$ with $\rho_i\geq0$. A rounded boundary collar gives the cofinal end family $L\times(0,\epsilon)$ and all bonding maps are homotopy equivalences.

For the nonarchimedean assertion, a point of $X^{\beth}$ is a seminorm with a center on $X$. A point of $W(X,D)$ has its generic point in $U$ and its center in $D$. Properness gives a unique lift of that center to $X'$, and separatedness makes the lift unique. Both tubes are thereby identified with the same analytic domain of $U^{\rm an}$, including its induced analytic structure. This is the valuative proof of Ben-Bassat--Temkin's comparison; see \cite[\S3.3.2 and Proposition~4.5.1]{BBT2013}. For $\mathbb G_m$ the two explicit domains are $0<|t|<1$ and the disjoint union of that domain with $1<|t|<\infty$. The natural map of the first to the second is the inclusion of one connected component. These comparisons prove the corrections to the added claim about cote~149, scan~133, and verify LM5 F.40--41 in the stated meanings.
\end{proof}

\begin{result}{D.4}{Proved}{The two requested ends and their comparison maps}
For the tripod, the archimedean end is three circles, the analytic open curve has homotopy type $S^1\vee S^1$, and the underlying trivially valued punctured tube is three open intervals. For $(\overline M_{0,5},D)$ the integral homology is free, with Betti numbers
\[
\begin{array}{c|rrrr}
 &b_0&b_1&b_2&b_3\\\hline
 M_{0,5}^{\rm an}&1&5&6&0\\
 L_{0,5}&1&11&11&1\\
 W(\overline M_{0,5},D)&1&6&0&0.
\end{array}
\]
The first comparison in D.1 sends every peripheral circle to its meridian. The nonarchimedean skeleton map sends a valuation to the normalized vector of its boundary orders, and retracts its underlying space onto the Petersen graph.
\end{result}
\begin{proof}
The tripod calculation follows by taking three punctured disks and a pair-of-pants core. A seminorm on a punctured formal disk over a trivially valued algebraically closed field is determined by $r=|t|\in(0,1)$: every unit has norm one, and a series of order $m$ has norm $r^m$.

Normalize three labels of $M_{0,5}$ to $0,1,\infty$. Projection $(x,y)\mapsto x$ from $\{x,y\in\mathbb C\setminus\{0,1\}:x\ne y\}$ is a bundle with base a twice-punctured plane and fibre a thrice-punctured plane. The monodromy fixes the oriented puncture meridians in first homology. The integral Serre spectral sequence has only two rows and columns and no possible differential; its ranks are the coefficients of $(1+2t)(1+3t)$.

The ten boundary divisors are indexed by two-element subsets of five labels; two meet exactly when the subsets are disjoint. They meet transversely once, and no three meet. Each is a rational curve of self-intersection $-1$ in the degree-five del Pezzo model. A regular neighborhood $N$ retracts to this union, so $H_1(N)=\mathbb Z^6$, $H_2(N)=\mathbb Z^{10}$, and $H_3(N)=0$. The intersection matrix is $Q=-I+A$, where $A$ is the Petersen adjacency matrix. The exact integral computation in \texttt{code/d\_boundary\_checks.py} gives Smith form with five unit diagonal entries and five zero entries. Poincar\'e--Lefschetz duality and the exact sequence of $(N,\partial N)$ then give $H_1(L)=\mathbb Z^{11}$ and $H_2(L)=\mathbb Z^{11}$; connectedness and orientation give $H_0=H_3=\mathbb Z$. This independently verifies LM5 F.42's archimedean row.

The local skeletal coordinates of the nonarchimedean tube are $a_i=-\log|z_i|\geq0$, not all zero. Normalization by $\sum a_i$ defines the comparison to the boundary complex. The toroidal retractions glue under boundary-coordinate changes and give the deformation retraction proved in Thuillier's theorem \cite[Theorem~3.25 and Proposition~4.5]{Thuillier2007}. The Petersen graph has ten vertices, fifteen edges and first Betti number six, proving the last row.
\end{proof}

\subsection{Rank selection and the levels}

\begin{result}{D.5}{Proved}{The exact link calculation and a uniform rank-selection bound}
For a nonempty multicurve $\sigma$ let $v_0(\sigma)$ be the number of rational components after cutting. Then its upper link has homotopy type a wedge of spheres of dimension
\[
 H(\sigma)=2g-3+n-v_0(\sigma).
\]
A pair-of-pants factor is interpreted as the join unit $S^{-1}$. For $g\geq1$ and $0\leq k<d$, $|P_k|$ is at least
\[
 \left(\min\left\{h(g,n),\left\lfloor\frac{2(k+1)}3\right\rfloor\right\}-1\right)\text{-connected}.
\]
For $g=0$ it is $(k-1)$-connected. In every genus $P_1$ is connected and $P_2$ is simply connected, with the empty-multicurve convention at $k\geq d$.
\end{result}
\begin{proof}
Write the cut components as $S_{g_i,n_i}$. Every component has $n_i>0$ and their complexes join to the link. The sphere dimension of their join plus one is the sum of the individual sphere dimensions plus one. With $v_+$ positive-genus vertices, $v=v_0+v_+$, and $q=|\sigma|$, the identities
\[
 \sum n_i=n+2q,\qquad \sum g_i=g-q+v-1
\]
give
\[
 \sum_{g_i>0}(2g_i-2+n_i)+\sum_{g_i=0}(n_i-3)-1
 =2g-3+n-v_0.
\]
This also directly types the join comparison; it is obtained by restricting disjoint curves to the cut components and taking their union for the inverse map.

Let $r=\sum(3g_i-3+n_i)=d-q$ be the residual complexity. A positive-genus factor of complexity $r_i>0$ has $g_i\leq\lfloor(r_i+2)/3\rfloor$, hence sphere dimension plus one at least $\lfloor2r_i/3\rfloor+1$. A rational factor contributes exactly $r_i$. For each positive integer $r_i$ these quantities are at least $2r_i/3+1/3$. Summing and taking the integer part proves $H(\sigma)\geq\lfloor2r/3\rfloor$ whenever $r>0$; this also holds when all factors are rational, since then $H=r-1\geq\lfloor2r/3\rfloor$ for $r\geq3$. The only smaller values $r=1,2$ are checked directly and give $H=0,1$.

Delete barycentric vertices of sizes $q=1,2,\ldots,d-k-1$ in that order. At deletion the remaining link is the upper link above and $r\geq k+1$. For $k\geq2$ the lower bound just proved applies; for $k=0,1$ use nonemptiness and the direct small-rank checks. Coning a $c$-connected link is a $(c+1)$-connected attachment and preserves homotopy groups through degree $c$. The initial curve complex is $(h-1)$-connected. Vertices of equal rank are incomparable, so all cones at one rank attach independently; weak CW realizations justify the infinite coproduct. This yields the asserted estimate. In genus zero every link has dimension $d-q-1$ and connectivity $d-q-2$, giving the sharper $(k-1)$ bound. For $k=1,2$ the displayed estimate gives respectively $0,1$; the exceptional low complexities are cones.
\end{proof}

\begin{result}{D.6}{Refuted}{The original Lead 4 and the four requested computations}
The original bound $(\min(k,2g-2+n)-1)$ is false: at $(g,n,k)=(2,1,3)$ the complex is a nontrivial wedge of two-spheres. The complete low-type table, with $\ast$ denoting a contractible cone and $\bigvee_\infty$ a countable wedge, is
\[
\begin{array}{c|c|c|c|c}
(g,n)&d&P_1&P_2&P_3\\\hline
(1,2)&2&\bigvee_\infty S^1&\ast&\ast\\
(1,3)&3&\bigvee_\infty S^1&\bigvee_\infty S^2&\ast\\
(2,0)&3&\bigvee_\infty S^1&\bigvee_\infty S^2&\ast\\
(2,1)&4&\bigvee_\infty S^1&\bigvee_\infty S^2&\bigvee_\infty S^2.
\end{array}
\]
All $P_0$ are nonempty discrete sets; $P_k$ is contractible for $k\geq d$. The corrected proposed bound with $h(g,n)$ in place of $2g-2+n$ holds in all four rows.
\end{result}
\begin{proof}
At $k=d-1$ the poset is all nonempty faces of the curve complex, so barycentric realization and the opened sphere theorem give the wedge entries. For $k<d-1$, the dimension of $P_k$ is $k$ and D.5 gives connectedness at $k=1$ and simple connectedness at $k=2$. Each connected countable graph is a wedge of circles. For $(1,3)$ and $(2,0)$ cut one nonseparating curve; its complement is respectively $(0,5)$ and $(1,2)$. The proper cofaces of that curve embed the barycentric subdivision of its curve complex as a subgraph of $P_1$. For $(2,1)$ use two nonseparating curves with connected complement $(0,5)$. These subgraphs have countably infinite first Betti number, and inclusion of a subgraph injects first homology, proving all circle multiplicities. For $P_2$ at $(2,1)$ cut one nonseparating curve. Its complement has type $(1,3)$, so its proper-coface subcomplex has countably infinite second homology. Inclusion into the two-dimensional $P_2$ injects second homology, because no three-chains exist. A simply connected two-dimensional CW complex is a wedge of two-spheres: choose a basis of its free second cellular homology, realize it by spheres using Hurewicz, and apply the homological Whitehead theorem. Countability then gives exactly the displayed multiplicity. At $(2,1,3)$ the second homology is a countably generated free abelian group, so the original asserted two-connectedness fails. This is an infinite topological counterexample; the finite calculations supplied with this section verify the graph types and rank data, not a replacement of the infinite curve complex by its finite quotient.
\end{proof}

\begin{result}{D.7}{Refuted; exact replacement proved}{The corrected Lead 4 also fails, with an explicit retraction}
For every $g\geq1$, $n\geq2$, put $d=3g-3+n$ and $h=2g-3+n$. Then $P_{d-2}$ is exactly $(h-2)$-connected: its first possible nonzero reduced homology is in degree $h-1$, and that group contains a direct summand
\[
 \widetilde H_{h-1}(\mathcal C_{g,n-1};\mathbb Z)
       \cong\bigoplus_{\mathbb N}\mathbb Z.
\]
In particular, at $(g,n,k)=(2,2,3)$ one has $d=5$, $h=3$, and $P_3$ is simply connected with nonzero second homology. This refutes the corrected proposed $(\min(k,h)-1)$ bound. More generally the same counterexample family refutes it for every $g\geq2,n\geq2$ at $k=d-2$. D.5 is the uniform all-genus bound proved here.
\end{result}
\begin{proof}
Label the punctures $p_1,\ldots,p_n$ and forget $p_n$. An essential curve becomes inessential or peripheral precisely when it bounds a disk containing $p_n$ and exactly one other puncture. No two distinct such curves can be disjoint. Indeed, their disks both contain $p_n$; disjoint boundary curves force the disks to be nested. If the other puncture is different, the outer disk contains at least three punctures and is not of the required kind. If it is the same, the annulus between the two boundaries contains no puncture, so the two curves are isotopic. Both alternatives exclude two distinct compatible vertices of this kind.

Every element $\sigma$ of $P_{d-2}$ has at least two curves. Delete any curve just described, fill in $p_n$, and identify any remaining curves which become isotopic. At least one essential curve remains, by the preceding paragraph. This defines a map
\[
 r:P_{d-2}(S_{g,n})\longrightarrow
       \operatorname{Face}^{+}(\mathcal C_{g,n-1}).
\]
On a curve which survives it is the actual filling map on its isotopy class, and on a killed curve it is deletion. Images of nested multicurves are nested even when duplicates merge, so $r$ is an order map and induces a continuous map on their order-complex realizations.

Fix a curve $b$ bounding a disk containing exactly $p_1,p_n$, and let $Y$ be its complementary component. Regard the boundary of $Y$ as a puncture; $Y$ has type $(g,n-1)$. Gluing back the disk after filling $p_n$ adds an annular collar terminating at the puncture $p_1$. Its inclusion induces an isomorphism
\[
 J:\mathcal C(Y)\xrightarrow{\sim}\mathcal C_{g,n-1},
\]
whose inverse restricts a curve after moving it out of that annular collar. This specifies the identification used below. For each nonempty multicurve $\tau$ in $Y$ set
\[
 i(\tau)=\tau\cup\{b\}.
\]
All its curves are essential and pairwise disjoint on $S_{g,n}$, and it has at least two members, so $i$ is an order map into $P_{d-2}$. The only deleted member of $i(\tau)$ is $b$; filling does not identify distinct curves of $Y$, by the collar identification. Consequently
\[
 r\circ i=J
\]
as order maps. Composing $r$ with $J^{-1}$ gives a genuine retraction of realizations onto the barycentric curve complex of $Y$. It follows on integral homology that $i_*$ is a split injection. Harer's sphere dimension for $Y$ is
$2g-3+(n-1)=h-1$, proving the stated nonzero direct summand.

It remains to prove the lower connectivity, so that the result is exact. Pass from the full nonempty barycentric curve complex to $P_{d-2}$ by deleting just the vertices corresponding to one-curve multicurves. The cut surface of a single essential curve has at most one rational component. By D.5 its upper link is a wedge of spheres of dimension either $h$ or $h-1$. Its connectivity is therefore at least $h-2$. The rank-deletion argument, starting with the $(h-1)$-connected full curve complex, proves that the remaining space is $(h-2)$-connected. At $h=1$ the assertion is simply nonemptiness and follows directly from the section $i$. Thus $P_{d-2}$ has exactly that connectivity.

Finally, $d-2-h=g-2$. For $g\geq2$ the proposed lower bound would be $h-1$, one higher than the exact value just proved. At $(2,2)$ the section embeds $\mathcal C_{2,1}\simeq\bigvee_{\mathbb N}S^2$ as a retract of $P_3$, giving the first displayed counterexample with all maps specified.
\end{proof}

The finite graph-type calculation in \path{code/d_rank_quotients.py} was also run. In particular its $(2,2),k=3$ quotient has zero reduced mod-two homology, while the actual complex just proved has nonzero integral second homology. The explicit quotient therefore does not detect this class. This is a verified instance of why the finite type-poset computation alone cannot settle the rank-selected question.


\subsection{Checking the projective threefold comparison}

In this subsection a \emph{good pair} is a smooth complex projective variety with reduced normal-crossings boundary, every connected open stratum, including $U$, aspherical, and every full normal-meridian lattice injecting into $\pi_1(U)$. The marking by $U$ is part of the data. The comparison below concerns these algebraically marked completions. It makes no assertion about arbitrary birational maps that move $U$.

\begin{result}{D.8}{Proved}{Good projective pairs are Mori hyperbolic and their comparisons are log crepant}
A good projective pair has no nonconstant algebraic map $\mathbb A^1\to S$ into any open stratum. For a simple normal-crossings boundary, $K_X+D$ is nef. For two such completions and a common smooth normal-crossings resolution $p:Z\to X$, $q:Z\to Y$, write $B=Z\setminus U$. Then
\[
 K_Z+B=p^*(K_X+D)+A_p=q^*(K_Y+E)+A_q,
 \qquad A_p=A_q\geq0.
\]
The common divisor is exceptional for both projections. A proper birational morphism between good smooth projective simple normal-crossings pairs, equal to the identity on $U$, is toroidal and log \'etale.
\end{result}
\begin{proof}
A map $f:\mathbb A^1\to S$ extends to $\bar f:\mathbb P^1\to X$ by properness. If infinity remains in $S$, asphericity makes the sphere nullhomotopic, whereas the integral of the pullback of a projective K\"ahler form is positive for a nonconstant holomorphic sphere. If infinity enters a deeper stratum $T$, trivialize the pulled-back old normal lines on the contractible source $\mathbb A^1$. Push the map a small nonzero distance in these directions into $U$, choosing the radius to tend to zero faster than every newly vanishing boundary coordinate near infinity. The boundary of a large disk then gives a nullhomotopic loop in $U$. In a normal-crossings chart at $\bar f(\infty)$ its newly acquired meridian coordinates are the nonzero contact orders of $\bar f$, with a common orientation sign. The old normal coordinates do not alter those new coordinates. This contradicts injection of the full lattice at $T$. Thus no such map exists. The open strata are exactly the strata in Svaldi's Theorem~1.1 \cite{Svaldi2019}; its hypotheses are smoothness, projectivity, reduced SNC boundary and the just-proved absence of these affine lines. Its conclusion gives nefness.

The logarithmic differential maps have effective determinant divisors $A_p,A_q$ supported on exceptional prime divisors. Put
\[
 R=p^*(K_X+D)-q^*(K_Y+E)=A_q-A_p.
\]
The divisor $-R$ is $p$-nef and $p_*R=p_*A_q$ is effective. The negativity lemma therefore gives $R\geq0$. Interchanging $p$ and $q$ gives $-R\geq0$, hence equality. This proves the claimed identity without asserting that the common resolution is good.

For an actual birational morphism $f:Y\to X$, the effective logarithmic Jacobian $A$ in $K_Y+E=f^*(K_X+D)+A$ is exceptional and $f$-nef. The same negativity lemma gives $A=0$. Thus the logarithmic differential matrix is everywhere invertible. In local boundary coordinates, $f^*x_i=u_i\prod_j y_j^{a_{ij}}$. The exponent matrix has full column rank; after choosing a nonsingular square minor and taking \'etale roots of units, absorb the corresponding units into the $y_j$. The remaining boundary-unit functions together with the ordinary target coordinates have independent ordinary differentials. They provide the invertible coordinates of a toric chart. Hence the morphism is locally an \'etale pullback of a monomial toric morphism; the finite lattice cokernel is invertible in characteristic zero. This is the explicit chart proof of log \'etaleness and toroidality.
\end{proof}

\begin{result}{D.9}{Proved}{Theorem R for projective threefolds, with the factorization reference corrected}
Two good smooth projective normal-crossings completions $(X,D)$ and $(Y,E)$ of the same smooth complex threefold $U$ admit a common good smooth projective completion $X_s$ such that $X_s\to X$ is a sequence of blow-ups of smooth closures of strata and $X_s\to Y$ is toroidal. They are joined over $U$ by a finite sequence of smooth stratal blow-ups and blow-downs through good projective pairs. This confirms Theorem~R of the attached review in its projective threefold scope.
\end{result}
\begin{proof}
First blow up the closures of the self-intersection strata in increasing dimension, obtaining simple normal-crossings boundaries. Locally these are barycentric subdivisions of the normal orthants. Every new open stratum is a torus bundle over an old open stratum, hence is aspherical by the homotopy exact sequence. Its normal-meridian lattice maps by a linearly independent set of primitive fan vectors into the old lattice, hence still injects into $\pi_1(U)$. This proves goodness at every step, and also proves the same preservation for the further stratal blow-ups below.

By D.8 every component $E_j$ of $E$ has log discrepancy zero over $(X,D)$. For a smooth SNC pair, the divisorial valuations of log discrepancy zero are exactly those produced by successive blow-ups of strata. One way to see the equality criterion is to pull back the logarithmic volume form: a center lying properly inside an open stratum supplies at least one ordinary coordinate and gives strictly positive discrepancy. A sequence following the center of the valuation therefore has discrepancy zero only at corner centers. In a normal orthant this is precisely a primitive monomial valuation. This criterion is also stated after Definition~2.7 in Svaldi's opened text, by reduction to the SNC discrepancy calculation. A primitive integral ray can be extracted by smooth star subdivisions. Write its coordinates in a smooth cone as $(a_1,\ldots,a_r)$ and blow up the face spanned by the basis vectors with positive coefficients. In the new cone indexed by a least positive coefficient $a_i$, the new coordinates are $a_i$ on the sum ray and $a_j-a_i$ on the other positive basis vectors. Their sum strictly decreases while at least two coefficients are positive. The process terminates at a coordinate ray, whose final coefficient is one by primitivity. Finitely many such extractions give a good smooth projective $X'\to X$ on which every prime divisor of $Y$ occurs.

Choose a very ample line bundle $H$ on $Y$. Transport its complete linear system to $X'$ and remove the common divisorial part. Let $I$ be the remaining base ideal and $L$ its line bundle. Its cosupport lies in $D'=X'\setminus U$ and has codimension at least two. Let
\[
 p:Z=\operatorname{Nor}(\operatorname{Bl}_I X')\longrightarrow X',
 \qquad q:Z\longrightarrow Y
\]
be the normalized graph. Writing $I\mathcal O_Z=\mathcal O_Z(-F)$, we have $q^*H=p^*L-F$ and $\mathcal O_Z(-F)$ is $p$-ample. Consequently $q$ is finite on every fibre of $p$: on such a fibre, the pullback of $H$ is ample, and a positive-dimensional fibre of $q$ would contain a curve on which that line bundle has degree zero.

Let $A$ be a $p$-exceptional prime divisor, equivalently a Rees divisor. It is absent as a divisor of $X'$, and therefore absent as a divisor of $Y$ by our extraction. Thus $\dim q(A)\leq1$. Its center on $X'$ cannot be a point, since then the two-dimensional $A$ would be mapped finitely by $q$ into a curve or a point. Hence its center is a curve $B$, and $q(A)$ is a curve $R$.

We now justify the crucial assertion that $B$ is a stratum curve and that the valuation is monomial. Choose a general point $b\in B$ avoiding every other exceptional center that does not dominate $B$, and a smooth surface slice transverse to $B$. Resolve the graph by principalizing $I$ on a smooth model. On the slice the resolution is a succession of point blow-ups of a smooth surface. Its exceptional divisor has a tree as dual graph, and its components are rational curves. The strict transforms of the boundary branches through $b$ give either one or two marked arrows. Each branch gives one arrow because its smooth curve germ has a single strict transform meeting the exceptional fibre in a single point. There are at most two branches at the generic point of a curve in a smooth SNC threefold.

Retain the vertices corresponding to Rees components, together with the marked arrows, and let $T_0$ be their minimal connecting subtree. A Rees component maps finitely to a rational curve $R_i\subset Y$. Let $S_i$ be the open stratum containing its generic point. The normalization of $R_i$ has at least two points outside $S_i$, by the absence of maps $\mathbb A^1\to S_i$ proved in D.8 (zero missing points also gives such a map). The same holds on its source rational component, because a finite surjection has a nonempty fibre over each of those two distinct points.

Every such boundary point comes from a boundary component $E_j$ not containing $R_i$. Its inverse image is an effective Cartier divisor on the normal graph $Z$, and therefore has pure codimension one. Here is the precise purity argument needed. Any supporting prime exceptional over $X'$ is a Rees divisor. Its center is a curve by the preceding paragraph; since $b$ avoids every other such center, that curve must be $B$. A supporting prime which is not exceptional is the strict transform of a prime divisor of $X'$. The map is the identity on $U$, so that prime lies in $D'$; it meets the general fibre over $b$ only when it contains $B$. These are exactly the one or two marked branches. The chosen Rees component is not itself in this support, since $E_j$ does not contain its image $R_i$.

On the resolved transverse surface, every exceptional component not corresponding to a Rees divisor contracts to a point of $Z$: a component mapping onto a curve of the general fibre would sweep out an exceptional prime of $Z$ dominating $B$, and would itself be a Rees component. Its connected non-Rees subtrees likewise contract. Thus all arrows and Rees components in any one direction of the resolved tree meet the chosen component at one point after those intervening subtrees are collapsed. A branch containing no arrow and no Rees component cannot supply an extra boundary point, because it supplies no supporting prime of the Cartier pullback on $Z$. It follows that every Rees vertex has at least two incident directions in $T_0$.

A finite tree with a retained vertex and only one marked arrow has an unmarked leaf. This contradicts the preceding assertion. Thus there are two arrows, so $B$ is an irreducible component of the intersection of two boundary divisors. Moreover $T_0$ is the path between the two arrows: a branch off this path would end at a Rees leaf. In the resolution of two transverse smooth curve germs, the vertices on that path are exactly those obtained by successive blow-ups of corner intersection points. Induction verifies this: a corner blow-up subdivides the path, while a blow-up at any other point adds a branch off it. Consequently every Rees valuation is a monomial valuation in the two boundary parameters. This also treats a Rees divisor whose generic fibre initially has several geometric components: the path is linearly ordered by its two labelled boundary arrows, and each corner component has its distinct primitive pair of weights. Monodromy preserves the two arrows and those weights, so it cannot permute different path positions. The valuation is therefore the indicated primitive monomial valuation over the function field of $B$.

For every $m\geq1$, the integral closure of $I^m$ is the intersection of the valuation ideals $v(f)\geq m v(I)$ over these finitely many Rees valuations. We use the existence and normalized-blow-up construction in Huneke--Swanson, Definition~10.1.1 and Theorem~10.2.2 \cite{HunekeSwanson2006}, both opened. Each inequality is monomial in boundary parameters, so each integral closure, and in particular $\overline I$, is a toroidal ideal sheaf. This includes points where several center curves meet. In a triple-crossing chart $x,y,z$, a Rees valuation over $x=y=0$ has weights $(a,b,0)$, and its valuation ideal is generated by the monomials $x^i y^j z^k$ with $ai+bj\geq m v(I)$. Valuations over $x=z=0$ or $y=z=0$ give the other two kinds of monomial inequality. At the special point the local integral closure is the intersection of all these valuation ideals. A finite intersection of monomial ideals is monomial: its generators are the least common multiples of choices of generators. The same argument applies to every subcollection of the stratum curves. An additional isolated integral-closure condition would require a Rees valuation centered at that point, already excluded by finiteness on $p$-fibres. Thus the normalized blow-up is toroidal everywhere, including the special triple points.

Principalize this monomial ideal by smooth stratal blow-ups. Canonical principalization, in the form recalled in AKMW \S1.2.3 \cite{AKMW2002}, is invariant under local torus actions and \'etale charts. Applied to a monomial ideal its smooth centers are unions of torus orbits closed in the chart, hence closures of strata; the local procedures agree on overlaps. We obtain a good $X_s\to X'$ and a morphism $X_s\to Y$. D.8 makes the latter toroidal. The corrected toroidal weak factorization theorem, AKMW Theorem~2.7.1 and Matsuki's erratum \S2 \cite{Matsuki2000}, factors it into smooth stratal blow-ups and blow-downs.

We use its relative morphism case. For the projective toroidal morphism $X_s\to Y$, the cobordism has a fan map to the fan of $Y$; its subdivisions preserve this map, and every elementary quotient fan still maps to that endpoint fan. These are the explicit fan maps in the proof of AKMW Theorem~2.7.1. Thus every intermediate variety is projective over $Y$ and its fan subdivides the fan of $Y$. This relative assertion proves goodness through the blow-down steps: every intermediate open stratum is a torus bundle over an open stratum of $Y$, and its meridian lattice injects, by the independent primitive generators of its cone, into the normal lattice of that endpoint stratum. Asphericity and meridian injectivity follow directly from goodness of $Y$, independently of the direction of the adjacent blow-up arrow. This gives the entire stated zigzag.

The proof uses only weak factorization. The algorithm for strong toroidal factorization in the earlier Morelli/AMR citation was withdrawn in the opened erratum and is not used.
\end{proof}

\begin{result}{D.10}{Proved}{All one-step stratal modifications of $\overline M_{0,6}$}
There are exactly $210$ nontrivial labelled smooth stratal blow-ups of $\overline M_{0,6}$: $105$ along boundary curves and $105$ at boundary points. Under permutation of the six labels they form four orbits, of sizes $45,60,15,90$. There is no nontrivial one-step stratal blow-down from the standard pair to another smooth SNC pair. The associated maps of conical complexes are the star subdivisions at the chosen two- or three-dimensional cones, with new ray the sum of their primitive generators. Theorem D.9 consequently applies to every good smooth projective completion of $M_{0,6}$.
\end{result}
\begin{proof}
A stable six-legged tree has at most three edges. The exact enumeration in \path{code/d_boundary_checks.py}, using the independently generated graph list of Part~A, gives $25,105,105$ strata of codimensions one, two, three. Blowing up a divisor is the identity. For two-edge trees the end/middle leg distributions are $(2,2,2)$ and $(2,1,3)$, with $45$ and $60$ labellings. For three-edge trees the internal graph is either a three-armed star, with two legs at each endpoint, or a path with endpoint leg counts two and interior leg counts one. Their counts are $15$ and $90$. Smoothness of the corresponding star subdivisions follows because each replacement maximal-cone matrix has determinant $1$; all these determinants are checked exactly by the program.

A smooth stratal blow-down in dimension three must contract a boundary divisor that is either a projective plane over a point or a projective-line bundle over a smooth curve. A divisor of type $2|4$ is $\overline M_{0,5}$, a rational surface with second Betti number five. It is neither $\mathbb P^2$ nor a projective-line bundle over a curve: rationality forces that curve to have genus zero, and such a bundle has second Betti number two. A divisor of type $3|3$ is $\mathbb P^1\times\mathbb P^1$. Any contraction to a curve uses a ruling, but its boundary contains three sections of that ruling. Their images would be three distinct boundary branches containing the entire center curve. Three SNC branches in a smooth threefold meet in codimension three, so this target would fail the SNC condition. These exhaust the twenty-five divisors. Finally, at each stratum of the standard moduli compactification the cut-surface Teichm\"uller groups are aspherical and independent commuting Dehn twists inject; hence the pair is good and D.9 applies.
\end{proof}

\subsection{Completion of the level complexes}

Fix a finite-index pure subgroup $G\subset\Gamma_{g,n}$ so that the simplicial action has no inversions; taking a barycentric subdivision also removes inversions. For a simplicial $G$-complex $K$ with finitely many cell orbits, its \emph{full level completion} is the simplicial profinite set
\[
 K^{\mathrm{lev}}=\varprojlim_{H\leq_fG}K/H.
\]
The procongruence version uses only congruence subgroups $H$. The map between them is induced by restriction of the inverse system. These definitions fix the exact completions used below.

\begin{result}{D.11}{Proved}{Profinite simple connectivity at level two and above}
For every stable type and every $k\geq2$, the full level completion of $P_k$ has trivial profinite fundamental group. The completed $P_1$ is connected. For the full curve complex, the same simple-connectivity conclusion holds for $d\geq3$. Completing the discrete presentation from the complex of groups on $P_2$ gives the full profinite group $\widehat\Gamma_{g,n}$ as a profinite presentation, with generators furnished by strata of modular dimension at most one and additional gluing relations furnished in dimension two.
\end{result}
\begin{proof}
D.5 and the cone convention show that $P_k$ is simply connected for $k\geq2$. For any finite-index $H\subset G$, Armstrong's quotient theorem gives
\[
 \pi_1(P_k/H)=H/H_*,
\]
where $H_*$ is the normal subgroup generated by elements with fixed points. The theorem's hypotheses hold: the complex is simply connected, the action is simplicial and without inversions. The same theorem and derivation are opened in Kent \S7, Theorem~11 \cite{Kent2016}. Thus $\widehat H$ surjects onto $\widehat{\pi_1(P_k/H)}$. Inverse limits of these surjections are surjective by compactness. But $\varprojlim_{H\leq_fG}\widehat H=1$: each $\widehat H$ identifies with the closure of $H$ in $\widehat G$, and their intersection is the identity by the defining finite-quotient topology. The profinite fundamental group commutes with this inverse limit of finite simplicial sets, in the sense specified in Kent \S5. Hence the desired group is trivial. The proof for the curve complex uses its simple connectivity at $d\geq3$. Connectedness of completed $P_1$ follows from connectedness of all its finite-level quotients and the corresponding inverse-limit statement for $\pi_0$.

For the presentation assertion, choose a maximal tree, stabilizer lifts, and transporter elements in the discrete complex of groups. Its ordinary fundamental group is $\Gamma_{g,n}$ because $P_2$ is simply connected. The homomorphisms from this presentation to any finite group are exactly the compatible homomorphisms from the local stabilizers and the edge transporters satisfying the edge and triangle relations. By the universal property of completion, homomorphisms from each local group to a finite group extend uniquely and continuously to its full completion. The same finite-group-valued universal property therefore identifies the resulting profinite presentation with $\widehat\Gamma_{g,n}$. This proves the presentation as a profinite colimit; it does not silently assert injectivity of every completed local-group map. Such injectivity is a separate comparison supplied where needed by Part~E.
\end{proof}

\begin{result}{D.12}{Refuted}{A printed low-complexity range in the profinite literature}
For the unmodified disjointness complex, the full level completion of $\mathcal C_{0,5}$ has a nontrivial continuous $\mathbb Z/2$-valued first cohomology class, and hence nontrivial profinite fundamental group. Thus the bound $3g-3+n\geq2$ printed in Kent's Theorem~12 cannot hold for the complex defined there. The argument in D.11 proves the corrected range $3g-3+n\geq3$. The genus-zero case of the analogous printed range in Theorem~13 likewise requires this correction.
\end{result}
\begin{proof}
Place the five labelled points on a circle. Let $c_{12},c_{23},c_{34},c_{45},c_{51}$ be boundaries of small regular neighborhoods of the five consecutive arcs. Curves associated with disjoint pairs are disjoint. Therefore
\[
 c_{12},c_{34},c_{51},c_{23},c_{45},c_{12}
\]
is a closed pentagon in the actual infinite curve complex. The quotient by the pure mapping class group maps a curve to the two-element side of its partition of the labels. It is the Petersen graph: edges join disjoint two-element subsets. The pentagon above maps to an embedded pentagon in that graph. On the Petersen graph take the $\mathbb Z/2$ cellular one-cochain equal to one on the edge $\{12,34\}$ and zero on all other edges. Every such cochain is a cocycle, since the graph has no two-cells, and its evaluation on this pentagon is one.

For every finite-index subgroup $H$ of the pure mapping class group, there are compatible maps
\[
 \mathcal C_{0,5}\longrightarrow\mathcal C_{0,5}/H
 \longrightarrow\mathcal C_{0,5}/\Gamma_{0,5}.
\]
The displayed closed pentagon supplies a closed lift in every intermediate graph. Thus the pulled-back cocycle remains nonzero at every level. It defines a nonzero continuous cocycle on the inverse-limit simplicial profinite graph: if it were a coboundary, a continuous zero-cochain with finite target would factor through one finite level, where the pentagon evaluation rules out a coboundary. This gives the asserted continuous character of the profinite fundamental group.

The finite Petersen graph, the lifted pair pattern, and the cocycle evaluation are checked in \texttt{code/d\_boundary\_checks.py}. No computation on a finite quotient is being used to prove that an arbitrary infinite loop survives: we supplied the actual closed lift. The source of the printed problem is visible in Kent's Theorem~10, which gives the correct Harer sphere dimensions but writes the complex dimension as $3g-3+n$ and infers simple connectivity one complexity too early. Armstrong's theorem requires that missing simple connectivity. The same sphere theorem gives $\mathcal C_{0,5}\simeq\bigvee_\infty S^1$ and D.11 supplies the corrected elementary argument.
\end{proof}

\begin{result}{D.13}{Known}{The proved procongruence range and the remaining scope}
For $n\geq2$, $2g-2+(n-1)>0$, and $3g-3+n\geq3$, the procongruence curve complex has trivial profinite fundamental group. This follows from Kent's Theorem~13 with its proof restricted to the actual simply connected range. The present proof does not establish the stronger higher-connectivity claims for arbitrary procongruence $P_k$, nor does it infer the congruence subgroup property from this one complex.
\end{result}
\begin{proof}
Put $S=S_{g,n-1}$ and $\dot S=S_{g,n}$. The first two inequalities make $S$ a smooth affine hyperbolic curve; the third makes $\mathcal C(\dot S)$ simply connected, which is the exact hypothesis required for the Armstrong sequence used in the proof. We spell out the comparison employed in Kent's \S8. Let $K$ be the congruence kernel, and let $\kappa$ range through finite-index congruence subgroups of the pure level subgroup. Armstrong's theorem, profinite right exactness, and inverse-limit exactness give
\[
 \varprojlim_\kappa\widehat{\kappa_*}\longrightarrow K
 \longrightarrow\Pi_1(\check{\mathcal C}(\dot S))\longrightarrow1.
\]
Choose a curve $\gamma$ bounding a pair of pants containing the added puncture and another puncture. Kent's Theorem~3 and Corollary~4 prove that the products $B_\kappa\kappa_\gamma$ are congruence and cofinal, where $B$ is the point-pushing subgroup and $B_\kappa=B\cap\kappa$. In the opened proof of Lemma~15, closure separation of the stabilizer implies
\[
 \bigcap_\kappa\overline{B_\kappa h_\gamma}
       =\overline{h_\gamma}
\]
for every fixed congruence subgroup $h$. Cofinality then places $K$ in the image of $\varprojlim\widehat{\kappa_\gamma}$, which lies in the image of $\varprojlim\widehat{\kappa_*}$. The displayed quotient is trivial. This is the exact step of the cited proof that transfers to our conventions. It uses the affine Birman sequence and the congruence cofinality proved in the source, not an assertion that all finite-index subgroups are congruence. D.12 explains why we have stated the necessary complexity restriction explicitly.
\end{proof}

\subsection{The twenty-three lambencies, computed as covers}

\begin{result}{D.14}{Proved}{All passports, cusp fields, Atkin--Lehner actions, and signatures}
For each lambency in Cheng--Duncan--Harvey's Table~2, the following tables give the complete permutation cover $X_0(N)\to X(1)$, its cusps, the action of the subgroup $S$ of exact divisors, and the signature of $\Gamma_0(N)+S$. The signature convention is $(0;e_2,e_3;c)$: genus zero, $e_i$ elliptic points of order $i$, and $c$ cusps. The genus $g_0$ in the first table is that of $X_0(N)$ before taking the $S$-quotient. A cycle expression $a^b$ means $b$ cycles of length $a$.

These are signatures of the effective modular orbifold. The homomorphism $\mathrm{SL}_2(\mathbb Z)\to\mathrm{PSL}_2(\mathbb Z)$ has kernel $\{\pm I\}$, which acts trivially on every dart. Pulling the effective quotient back through this central extension restores the moduli stack of elliptic curves: its generic inertia is $\mu_2$ and its elliptic inertia orders are $4$ and $6$. Thus the displayed finite cover and its unrigidified modular stack are related by that specified central gerbe.
\end{result}

{\small\begin{longtable}{r|r|r|l|l|l|l}
$\ell$ & $\mu$ & $g_0$ & $\sigma_2$ & $\sigma_3$ & cusp widths & $(e_2,e_3;c)$\\\hline
$2$ & 3 & 0 & $1 2$ & $3$ & $1 2$ & $(1,0;2)$\\
$3$ & 4 & 0 & $2^{2}$ & $1 3$ & $1 3$ & $(0,1;2)$\\
$4$ & 6 & 0 & $2^{3}$ & $3^{2}$ & $1^{2} 4$ & $(0,0;3)$\\
$5$ & 6 & 0 & $1^{2} 2^{2}$ & $3^{2}$ & $1 5$ & $(2,0;2)$\\
$6$ & 12 & 0 & $2^{6}$ & $3^{4}$ & $1 2 3 6$ & $(0,0;4)$\\
$7$ & 8 & 0 & $2^{4}$ & $1^{2} 3^{2}$ & $1 7$ & $(0,2;2)$\\
$8$ & 12 & 0 & $2^{6}$ & $3^{4}$ & $1^{2} 2 8$ & $(0,0;4)$\\
$9$ & 12 & 0 & $2^{6}$ & $3^{4}$ & $1^{3} 9$ & $(0,0;4)$\\
$10$ & 18 & 0 & $1^{2} 2^{8}$ & $3^{6}$ & $1 2 5 10$ & $(2,0;4)$\\
$12$ & 24 & 0 & $2^{12}$ & $3^{8}$ & $1^{2} 3^{2} 4 12$ & $(0,0;6)$\\
$13$ & 14 & 0 & $1^{2} 2^{6}$ & $1^{2} 3^{4}$ & $1 13$ & $(2,2;2)$\\
$16$ & 24 & 0 & $2^{12}$ & $3^{8}$ & $1^{4} 4 16$ & $(0,0;6)$\\
$18$ & 36 & 0 & $2^{18}$ & $3^{12}$ & $1^{3} 2^{3} 9 18$ & $(0,0;8)$\\
$25$ & 30 & 0 & $1^{2} 2^{14}$ & $3^{10}$ & $1^{5} 25$ & $(2,0;6)$\\
$6+3$ & 12 & 0 & $2^{6}$ & $3^{4}$ & $1 2 3 6$ & $(2,0;2)$\\
$10+5$ & 18 & 0 & $1^{2} 2^{8}$ & $3^{6}$ & $1 2 5 10$ & $(3,0;2)$\\
$14+7$ & 24 & 1 & $2^{12}$ & $3^{8}$ & $1 2 7 14$ & $(4,0;2)$\\
$18+9$ & 36 & 0 & $2^{18}$ & $3^{12}$ & $1^{3} 2^{3} 9 18$ & $(2,0;4)$\\
$22+11$ & 36 & 2 & $2^{18}$ & $3^{12}$ & $1 2 11 22$ & $(6,0;2)$\\
$30+15$ & 72 & 3 & $2^{36}$ & $3^{24}$ & $1 2 3 5 6 10 15 30$ & $(8,0;4)$\\
$46+23$ & 72 & 5 & $2^{36}$ & $3^{24}$ & $1 2 23 46$ & $(12,0;2)$\\
$12+4$ & 24 & 0 & $2^{12}$ & $3^{8}$ & $1^{2} 3^{2} 4 12$ & $(0,0;4)$\\
$30+6,10,15$ & 72 & 3 & $2^{36}$ & $3^{24}$ & $1 2 3 5 6 10 15 30$ & $(6,0;2)$\\
\end{longtable}
}

\begin{proof}
Use primitive row vectors $(x,y)$ modulo $N$, modulo multiplication by a unit. These are the cosets $\Gamma_0(N)\backslash\mathrm{SL}_2(\mathbb Z)$, by sending a matrix to its bottom row. Its stabilizer is exactly $\Gamma_0(N)$; lifting a primitive residue row to a primitive integral row and completing it to a determinant-one matrix proves surjectivity. The matrices
\[
 s=\begin{pmatrix}0&-1\\1&0\end{pmatrix},\quad
 r=\begin{pmatrix}0&-1\\1&1\end{pmatrix},\quad
 t=\begin{pmatrix}1&1\\0&1\end{pmatrix}
\]
act on the right, with $s^2=r^3=1$ and $sr=t$ in the projective action. Their cycles give the passport. Their numbers of fixed darts give $e_2,e_3$; the $t$-cycles give the cusp widths. Riemann--Hurwitz gives
\[
 g_0=1+\mu/12-e_2/4-e_3/3-c/2.
\]
The executable \texttt{d\_lambencies.py} constructs every primitive row, canonicalizes its unit orbit, and stores all permutations, rather than only their cycle partitions.

Write $f_d=\gcd(d,N/d)$. The cusp label $c_{d,a}$ denotes the class of $a'/d$, where $d\mid N$, $a\in(\mathbb Z/f_d)^\times$, and $a'\equiv a\pmod{f_d}$ is chosen coprime to $d$. For $f_d=1$ we write $a=1$. In particular $c_{1,1}=0$ and $c_{N,1}=\infty$. A reduced rational cusp $x/y$ has label
\[
 d=\gcd(y,N),\qquad a=x(y/d)\pmod{f_d}.
\]
For $\left(\begin{smallmatrix}A&B\\NC&D\end{smallmatrix}\right)\in\Gamma_0(N)$, its transformed numerator is $Ax$ modulo $f_d$, its denominator divided by $d$ is $D(y/d)$ modulo $f_d$, and $AD\equiv1$. This proves invariance. Conversely, completing the two primitive pairs to unimodular matrices and solving the lower-left congruence by B\'ezout gives equivalence exactly for equal $d,a$; equivalently this is the usual cusp criterion used in the opened modular-curve source \cite{RouseWebb2015}. Its width is obtained from the lower-left entry of $\gamma t^w\gamma^{-1}$: the least positive $w$ with $N\mid d^2w$ is
\[
 w_d=N/\gcd(d^2,N).
\]
The Galois action in \cite[Theorem~16]{RouseWebb2015} sends the numerator invariant to its product with a cyclotomic unit (or its inverse, depending on the chosen action convention). Since that action on $(\mathbb Z/f_d)^\times$ is free and transitive, the exact residue field is $\mathbb Q(\zeta_{f_d})$. This also supplies a direct derivation of the field, rather than merely a field over which the cusp becomes rational.

For an exact divisor $e$, choose integers $b,d_0$ with $ed_0-(N/e)b=1$. The matrix
\[
 W_e=\begin{pmatrix}e&b\\N&ed_0\end{pmatrix},\qquad\det W_e=e,
\]
acts on rational cusps; scaling it by $1/\sqrt e$ gives the standard Atkin--Lehner coset. Reduce the resulting fraction and apply the preceding cusp formula. The program checks all products $W_eW_f=W_{ef/\gcd(e,f)^2}$ on cusps. The denominator changes by $d\mapsto ed/\gcd(e,d)^2$, while the table retains the numerator action as well. The latter matters at $N=18$: $W_9$ exchanges the two cusps at denominator $3$ and the two at denominator $6$. The exact cusp fields on a quotient are fixed fields of the stabilizers of these finite Galois orbit sets. In every nontrivial $S$-quotient in this table all cusps are rational. This includes the two merged conjugate pairs for $18+9$.

D.15 below proves genus zero directly from the eta divisor and proves $S$-invariance. For the added cosets, the normalized trace is $\sqrt e(a+d)$. For $e\geq4$ an elliptic element must have trace zero, hence order two. In the sole $e=3$ case, $N=6$, the determinant equation $3ad-2bc=1$ makes $a,d$ odd, so again a trace of absolute value less than two must be zero. The identity coset has no order-three elliptic points for any of the nine nontrivial $S$-rows, as the explicit dart permutations verify. Thus only order two occurs. Hyperbolic area divided by $|S|$ gives
\[
 \frac{\mu}{6|S|}=-2+c+e_2/2,
\]
which proves every remaining signature entry.
\end{proof}

\begin{result}{D.15}{Proved}{The exact tripod comparison and the eta-product check}
The normalization of $X_0(N)\times_{X(1)}X(2)$ has one component of degree $\mu$ over $X(2)$ for odd $N$, and three components, each of degree $\mu/3$, for even $N$. They are defined over $\mathbb Q$. The component containing the standard cusp is isomorphic to $X_0(4N)$ for odd $N$ and to $X_0(2N)$ for even $N$; the other components differ by permutation of the three tripod punctures. For $4\mid N$ there is also the direct cover $X_0(N)\to X_0(4)$ of degree $\mu/6$.

For all twenty-three rows the exact modular function $T^X=1/\eta_{\pi^X}$ is $S$-invariant and descends to a $\mathbb Q$-coordinate on $X_0(N)/S$, with a unique simple pole and expansion $q^{-1}-d^X+O(q)$. The attached computation supplies the full integral cusp divisor and coefficients through $q^{30}$. The assertions of genus zero and invariance follow from those exact divisors and modular transformation criteria, without an inference from a truncated series.
\end{result}
\begin{proof}
The subgroup $\Gamma(2)$ is generated projectively by $t^2$ and $u_2=\left(\begin{smallmatrix}1&0\\-2&1\end{smallmatrix}\right)$. Its orbits on the dart set are the components of the normalized base change, and the three cusp permutations are $t^2,u_2,(t^2u_2)^{-1}$. For odd $N$, the Chinese remainder theorem makes the image of $\Gamma(2)$ act transitively on primitive rows modulo $N$. For even $N$, reduction modulo $2$ records three primitive row classes, and the same lifting argument makes each class an orbit. All three choices are defined over $\mathbb Q$, since the full two-torsion labeling defining $X(2)$ is defined over $\mathbb Q$. Conjugation by $D=\operatorname{diag}(2,1)$ gives, by direct matrix multiplication,
\[
 D\Gamma_0(4N)D^{-1}=\Gamma(2)\cap\Gamma_0(N)\quad(N\text{ odd}),
\]
\[
 D\Gamma_0(2N)D^{-1}=\Gamma(2)\cap\Gamma_0(N)\quad(N\text{ even}).
\]
This conjugation is the algebraic two-isogeny map with the induced two-torsion labeling, so defines the asserted comparison over $\mathbb Q$. The other components are obtained from the rational $S_3$-action on $X(2)$. For the direct cover use $\Gamma_0(N)\subset\Gamma_0(4)$ and the generators $t,u_4=\left(\begin{smallmatrix}1&0\\-4&1\end{smallmatrix}\right)$. Restricting their dart permutations to the orbit of $(0,1)$ computes that cover and its three-point passport. The direct passports for $N=4,8,12,16$ are respectively
\[
 (1,1,1),\quad(1^2,2,2),\quad(1\,3,1\,3,1\,3),\quad(1^4,4,4).
\]
All component permutations, not just their lengths, are in the JSON certificate.

Let $r_\delta$ be the eta exponents of $T^X$, so $r_\delta=-k_\delta$ for the displayed Frame shape. The opened Newman criterion and Ligozat formula in \cite[p.~2]{RouseWebb2015} apply: the program checks
\[
 \sum r_\delta=0,\quad\sum\delta r_\delta\equiv
 \sum(N/\delta)r_\delta\equiv0\pmod{24},\quad
 \prod\delta^{r_\delta}\in(\mathbb Q^\times)^2.
\]
Thus $T^X$ is a weight-zero function with trivial character on $\Gamma_0(N)$ and no zeros or poles in the upper half-plane. At $c_{d,a}$ its order is the exact rational sum
\[
 \operatorname{ord}_{c_{d,a}}T^X=
 \frac N{24}\sum_{\delta\mid N}
 \frac{\gcd(d,\delta)^2r_\delta}{\gcd(d,N/d)d\delta}.
\]
Every value is integral. For every row the negative part consists exactly of the $|S|$ cusps in the orbit of infinity, each with order $-1$; the whole divisor is invariant under the computed $S$-permutations. Consequently $(T^X\circ W_e)/T^X$ is a nonzero constant.

Here are exact constant-one certificates. For each D-type $N=2e$, the matrix
\[
 \begin{pmatrix}e&-(e+1)/2\\2e&-e\end{pmatrix}
\]
has determinant $e$ and trace zero. Its fixed point lies in the upper half-plane, where $T^X$ is finite and nonzero, forcing the constant to be one. At $(N,e)=(30,6)$ choose the trace-zero matrix $(12,-5;30,-12)$, and at $(30,15)$ choose $(15,-8;30,-15)$; their determinant checks are $6$ and $15$. These generate the group containing $W_{10}=W_6W_{15}$. For $(N,e)=(12,4)$, the cusp $c_{2,1}$ is fixed and has eta order zero, again forcing the constant to be one. These certificates are also checked by the program. Therefore $T^X$ descends to the quotient and has one simple pole there. A nonconstant meromorphic function of degree one identifies that smooth projective quotient with $\mathbb P^1$.

The $q$-product has leading order $-1$ and leading coefficient one; its next coefficient is $k_1=-d^X$, as checked against every root system of the opened CDH Table~2 \cite{CDH2014}. Rational $q$-expansions on the canonical rational modular curve and rationality of the Atkin--Lehner maps give the coordinate over $\mathbb Q$. This independently verifies the content of CDH Proposition~2.1 for the table used here.

The quotient maps retain a precise limitation on these covers. For every nontrivial $e$, the function $j\circ W_e$ has pole order $e$ at infinity, while $j$ has pole order one. Hence $j$ does not factor through $X_0(N)/\langle W_e\rangle$. In the special direct-cover row $12+4$, the map $X_0(12)\to X_0(4)$ sends infinity to infinity but sends $W_4(\infty)=c_{3,1}$ to the cusp zero. It therefore does not factor through that quotient either. The explicit covers above, together with the stated quotient maps, are the resulting tripod form; no descent of those particular maps is asserted.
\end{proof}

The next tables make each cusp field and every nontrivial $S$-action explicit. A set in the second subscript lists the distinct numerator invariants; it is not an unreduced rational fraction. The fields with $f=1,2$ are $\mathbb Q$. In the tripod table the other even-level components have the displayed passport with the three punctures permuted. The JSON and the full readable certificate \texttt{code/d\_lambencies.md} additionally list every quotient orbit and every individual eta order.

{\small\begin{longtable}{r|p{0.79\textwidth}}
$N$ & Cusps $c_{d,a}$, with $(w;f)$ meaning width $w$ and field $\mathbb Q(\zeta_f)$\\\hline
2 & $c_{1,\{1\}}:(2;1)$;\ $c_{2,\{1\}}:(1;1)$\\
3 & $c_{1,\{1\}}:(3;1)$;\ $c_{3,\{1\}}:(1;1)$\\
4 & $c_{1,\{1\}}:(4;1)$;\ $c_{2,\{1\}}:(1;2)$;\ $c_{4,\{1\}}:(1;1)$\\
5 & $c_{1,\{1\}}:(5;1)$;\ $c_{5,\{1\}}:(1;1)$\\
6 & $c_{1,\{1\}}:(6;1)$;\ $c_{2,\{1\}}:(3;1)$;\ $c_{3,\{1\}}:(2;1)$;\ $c_{6,\{1\}}:(1;1)$\\
7 & $c_{1,\{1\}}:(7;1)$;\ $c_{7,\{1\}}:(1;1)$\\
8 & $c_{1,\{1\}}:(8;1)$;\ $c_{2,\{1\}}:(2;2)$;\ $c_{4,\{1\}}:(1;2)$;\ $c_{8,\{1\}}:(1;1)$\\
9 & $c_{1,\{1\}}:(9;1)$;\ $c_{3,\{1,2\}}:(1;3)$;\ $c_{9,\{1\}}:(1;1)$\\
10 & $c_{1,\{1\}}:(10;1)$;\ $c_{2,\{1\}}:(5;1)$;\ $c_{5,\{1\}}:(2;1)$;\ $c_{10,\{1\}}:(1;1)$\\
12 & $c_{1,\{1\}}:(12;1)$;\ $c_{2,\{1\}}:(3;2)$;\ $c_{3,\{1\}}:(4;1)$;\ $c_{4,\{1\}}:(3;1)$;\ $c_{6,\{1\}}:(1;2)$;\ $c_{12,\{1\}}:(1;1)$\\
13 & $c_{1,\{1\}}:(13;1)$;\ $c_{13,\{1\}}:(1;1)$\\
14 & $c_{1,\{1\}}:(14;1)$;\ $c_{2,\{1\}}:(7;1)$;\ $c_{7,\{1\}}:(2;1)$;\ $c_{14,\{1\}}:(1;1)$\\
16 & $c_{1,\{1\}}:(16;1)$;\ $c_{2,\{1\}}:(4;2)$;\ $c_{4,\{1,3\}}:(1;4)$;\ $c_{8,\{1\}}:(1;2)$;\ $c_{16,\{1\}}:(1;1)$\\
18 & $c_{1,\{1\}}:(18;1)$;\ $c_{2,\{1\}}:(9;1)$;\ $c_{3,\{1,2\}}:(2;3)$;\ $c_{6,\{1,2\}}:(1;3)$;\ $c_{9,\{1\}}:(2;1)$;\ $c_{18,\{1\}}:(1;1)$\\
22 & $c_{1,\{1\}}:(22;1)$;\ $c_{2,\{1\}}:(11;1)$;\ $c_{11,\{1\}}:(2;1)$;\ $c_{22,\{1\}}:(1;1)$\\
25 & $c_{1,\{1\}}:(25;1)$;\ $c_{5,\{1,2,3,4\}}:(1;5)$;\ $c_{25,\{1\}}:(1;1)$\\
30 & $c_{1,\{1\}}:(30;1)$;\ $c_{2,\{1\}}:(15;1)$;\ $c_{3,\{1\}}:(10;1)$;\ $c_{5,\{1\}}:(6;1)$;\ $c_{6,\{1\}}:(5;1)$;\ $c_{10,\{1\}}:(3;1)$;\ $c_{15,\{1\}}:(2;1)$;\ $c_{30,\{1\}}:(1;1)$\\
46 & $c_{1,\{1\}}:(46;1)$;\ $c_{2,\{1\}}:(23;1)$;\ $c_{23,\{1\}}:(2;1)$;\ $c_{46,\{1\}}:(1;1)$\\
\end{longtable}
\begin{longtable}{r|p{0.78\textwidth}}
$\ell$ & Atkin--Lehner permutations; omitted cusp labels are fixed\\\hline
$6+3$ & $W_{3}=(c_{1,1}\ c_{3,1})(c_{2,1}\ c_{6,1})$\\
$10+5$ & $W_{5}=(c_{1,1}\ c_{5,1})(c_{2,1}\ c_{10,1})$\\
$14+7$ & $W_{7}=(c_{1,1}\ c_{7,1})(c_{2,1}\ c_{14,1})$\\
$18+9$ & $W_{9}=(c_{1,1}\ c_{9,1})(c_{2,1}\ c_{18,1})(c_{3,1}\ c_{3,2})(c_{6,1}\ c_{6,2})$\\
$22+11$ & $W_{11}=(c_{1,1}\ c_{11,1})(c_{2,1}\ c_{22,1})$\\
$30+15$ & $W_{15}=(c_{1,1}\ c_{15,1})(c_{2,1}\ c_{30,1})(c_{3,1}\ c_{5,1})(c_{6,1}\ c_{10,1})$\\
$46+23$ & $W_{23}=(c_{1,1}\ c_{23,1})(c_{2,1}\ c_{46,1})$\\
$12+4$ & $W_{4}=(c_{1,1}\ c_{4,1})(c_{3,1}\ c_{12,1})$\\
$30+6,10,15$ & $W_{6}=(c_{1,1}\ c_{6,1})(c_{2,1}\ c_{3,1})(c_{5,1}\ c_{30,1})(c_{10,1}\ c_{15,1})$\newline $W_{10}=(c_{1,1}\ c_{10,1})(c_{2,1}\ c_{5,1})(c_{3,1}\ c_{30,1})(c_{6,1}\ c_{15,1})$\newline $W_{15}=(c_{1,1}\ c_{15,1})(c_{2,1}\ c_{30,1})(c_{3,1}\ c_{5,1})(c_{6,1}\ c_{10,1})$\\
\end{longtable}
\begin{longtable}{r|r|r|l|l}
$N$ & components & degree & passport of one component over $X(2)$ & direct degree to $X_0(4)$\\\hline
2 & 3 & 1 & $(1,1,1)$ & --\\
3 & 1 & 4 & $(1\,3,1\,3,1\,3)$ & --\\
4 & 3 & 2 & $(1^{2},2,2)$ & 1\\
5 & 1 & 6 & $(1\,5,1\,5,1\,5)$ & --\\
6 & 3 & 4 & $(1\,3,1\,3,1\,3)$ & --\\
7 & 1 & 8 & $(1\,7,1\,7,1\,7)$ & --\\
8 & 3 & 4 & $(1^{4},4,4)$ & 2\\
9 & 1 & 12 & $(1^{3}\,9,1^{3}\,9,1^{3}\,9)$ & --\\
10 & 3 & 6 & $(1\,5,1\,5,1\,5)$ & --\\
12 & 3 & 8 & $(1^{2}\,3^{2},2\,6,2\,6)$ & 4\\
13 & 1 & 14 & $(1\,13,1\,13,1\,13)$ & --\\
14 & 3 & 8 & $(1\,7,1\,7,1\,7)$ & --\\
16 & 3 & 8 & $(1^{4}\,2^{2},8,8)$ & 4\\
18 & 3 & 12 & $(1^{3}\,9,1^{3}\,9,1^{3}\,9)$ & --\\
22 & 3 & 12 & $(1\,11,1\,11,1\,11)$ & --\\
25 & 1 & 30 & $(1^{5}\,25,1^{5}\,25,1^{5}\,25)$ & --\\
30 & 3 & 24 & $(1\,3\,5\,15,1\,3\,5\,15,1\,3\,5\,15)$ & --\\
46 & 3 & 24 & $(1\,23,1\,23,1\,23)$ & --\\
\end{longtable}
\begin{longtable}{r|l|r|r}
$\ell$ & $\pi^X$ (positive numerator, negative denominator) & $[q^0]T^X$ & $[q]T^X$\\\hline
$2$ & $\frac{2^{24}}{1^{24}}$ & -24 & 276\\
$3$ & $\frac{3^{12}}{1^{12}}$ & -12 & 54\\
$4$ & $\frac{4^{8}}{1^{8}}$ & -8 & 20\\
$5$ & $\frac{5^{6}}{1^{6}}$ & -6 & 9\\
$6$ & $\frac{2\,6^{5}}{1^{5}\,3}$ & -5 & 6\\
$7$ & $\frac{7^{4}}{1^{4}}$ & -4 & 2\\
$8$ & $\frac{2^{2}\,8^{4}}{1^{4}\,4^{2}}$ & -4 & 4\\
$9$ & $\frac{9^{3}}{1^{3}}$ & -3 & 0\\
$10$ & $\frac{2\,10^{3}}{1^{3}\,5}$ & -3 & 1\\
$12$ & $\frac{2^{2}\,3\,12^{3}}{1^{3}\,4\,6^{2}}$ & -3 & 2\\
$13$ & $\frac{13^{2}}{1^{2}}$ & -2 & -1\\
$16$ & $\frac{2\,16^{2}}{1^{2}\,8}$ & -2 & 0\\
$18$ & $\frac{2\,3\,18^{2}}{1^{2}\,6\,9}$ & -2 & 0\\
$25$ & $\frac{25}{1}$ & -1 & -1\\
$6+3$ & $\frac{2^{6}\,6^{6}}{1^{6}\,3^{6}}$ & -6 & 15\\
$10+5$ & $\frac{2^{4}\,10^{4}}{1^{4}\,5^{4}}$ & -4 & 6\\
$14+7$ & $\frac{2^{3}\,14^{3}}{1^{3}\,7^{3}}$ & -3 & 3\\
$18+9$ & $\frac{2^{3}\,3^{2}\,18^{3}}{1^{3}\,6^{2}\,9^{3}}$ & -3 & 3\\
$22+11$ & $\frac{2^{2}\,22^{2}}{1^{2}\,11^{2}}$ & -2 & 1\\
$30+15$ & $\frac{2^{2}\,3\,5\,30^{2}}{1^{2}\,6\,10\,15^{2}}$ & -2 & 1\\
$46+23$ & $\frac{2\,46}{1\,23}$ & -1 & 0\\
$12+4$ & $\frac{2^{4}\,3^{4}\,12^{4}}{1^{4}\,4^{4}\,6^{4}}$ & -4 & 6\\
$30+6,10,15$ & $\frac{2^{3}\,3^{3}\,5^{3}\,30^{3}}{1^{3}\,6^{3}\,10^{3}\,15^{3}}$ & -3 & 3\\
\end{longtable}
}


\subsection{The reconstruction, its stage, and the five link diagrams}

\begin{result}{D.16}{Proved}{Canonical reconstruction from the supplied face and collar diagram}
Let $\mathcal V$ be the real oriented blow-up of a smooth complex normal-crossings Deligne--Mumford pair $(\mathcal X,\mathcal D)$, with its full deployed face diagram, collars, and maps to the strata. Collapsing precisely the normal angular circles over the zero radial coordinates recovers $\mathcal X$ as a topological stack. Equivalently, the mapping-cylinder pushout
\[
 \widehat X=X^*\mathop{\amalg}_{\dot X} C(p)
\]
of cote~150, scans~32--35, performed successively on the deployed strata, gives this same stack. For the supplied data, the final isomorphism is canonical and independent of the order of a compatible induction.

For $\overline{\mathcal M}_{g,n}$, the fundamental groupoid is already that of the compactification after all open divisors have been attached to the open stratum. In the deepest-strata-first reconstruction starting from $\mathcal V$, it is already that groupoid after the first layer, consisting of all maximal-corner fibres. In both descriptions it is the connected simply connected groupoid. These are two explicitly related filtrations of the same final quotient.
\end{result}
\begin{proof}
On a uniformizing chart, write the boundary equation as $z_1\cdots z_q=0$. The blow-up and its blow-down are
\[
 (\rho_i,\theta_i,z_{q+1},\ldots,z_m)\longmapsto
 (\rho_i e^{\mathrm i\theta_i},z_{q+1},\ldots,z_m),\qquad \rho_i\geq0.
\]
Its fibres identify exactly the $\theta_i$ for which $\rho_i=0$. The resulting quotient chart is the original complex chart: polar coordinates give a continuous bijection from the compact local quotient to its Hausdorff image, and the same assertion on relatively compact subcharts gives a homeomorphism. The construction is equivariant for the finite chart group and for coordinate changes. Taking the quotient groupoids therefore gives the original topological stack.

For a closed attaching stratum $S$, let $p:L_S\to S$ be its deployed normal link. Its mapping cylinder is
\[
 C(p)=\bigl(L_S\times[0,1]\amalg S\bigr)/((\ell,0)\sim p(\ell)).
\]
The collar identifies $L_S\times\{1\}$ with the attaching boundary of the previously cut piece. Gluing this cylinder implements exactly the preceding angular quotient on the corresponding normal cone. Repeating on all faces recovers the common quotient of $\mathcal V$. The coordinate equivalence relations commute: collapsing coordinate $i$ and then $j$ identifies exactly the same pairs as collapsing $j$ and then $i$. More generally the universal property of the quotient by the union of these relations identifies every compatible iterated pushout with one fixed quotient. The induced isomorphisms are unique over that quotient and satisfy their cocycle law. This is the asserted canonical independence for the given face/collar data; it requires no claim that collars have a preferred choice.

For the first filtration, set $\mathcal X^{\leq1}=\mathcal M_{g,n}\cup\{\text{open boundary divisors}\}$. Its inclusion into $\overline{\mathcal M}_{g,n}$ induces an isomorphism on $\pi_1$: the omitted strata have complex codimension at least two, and a loop or a homotopy of loops can be moved off them by real transversality in the uniformizing charts. Van Kampen for attaching an open divisor kills its normal meridian. Every essential simple closed curve gives a boundary degeneration, so the resulting quotient is $\Gamma_{g,n}$ modulo the normal subgroup generated by all Dehn twists. These twists generate the pure mapping class group. In positive genus this is exactly Labru\`ere--Paris, Corollary~2.11(i), with its puncture-induction proof on pages~86--89 \cite{D_LabruereParis2001}; in genus zero it follows from the pure-braid generators given in Part~B. Thus this quotient is trivial.

For the second filtration, the inclusion of the interior into $\mathcal V$ is a homotopy equivalence by pushing inward along a collar, so $\pi_1(\mathcal V)=\Gamma_{g,n}$. Over a maximal boundary stratum the deployed face is a torus bundle over the finite automorphism stack of a pants graph. The map to that stack is surjective on $\pi_1$, with kernel the image of its normal torus lattice. Replacing it by its mapping cylinder therefore imposes only the relations killing that lattice. Every essential curve extends to a pants decomposition; the normal closure of the lattices over all maximal strata is consequently the subgroup generated by all Dehn twists. The first complete layer has trivial fundamental group. Subsequent angular collapses again have connected fibres and induce quotients on $\pi_1$, so the group stays trivial. These intermediate quotients need not themselves be smooth stacks; they are the topological pushouts specified by the reconstruction. The two filtrations map to the same final quotient just proved.
\end{proof}

\paragraph{The four reconstruction cases.}
For $(0,4)$, $\mathcal V$ is a pair of pants. The three face maps are $S^1\to\{\text{point}\}$, and their mapping cylinders are disks. The pushout is $\mathbb P^1$. Killing the three peripheral loops already gives its trivial fundamental group.

For $(0,5)$, the ten deployed codimension-one faces are the normal circle bundles over the cut versions of the ten boundary copies of $\mathbb P^1$; their open bases are tripods. Fifteen deployed corner fibres are two-tori, one for each edge of the Petersen incidence graph. Each of the ten normal line bundles has degree $-1$. On a corner the two collapse maps are $(e^{\mathrm i\theta_1},e^{\mathrm i\theta_2})\mapsto e^{\mathrm i\theta_1}$ and $(e^{\mathrm i\theta_1},e^{\mathrm i\theta_2})\mapsto e^{\mathrm i\theta_2}$, followed by the corresponding remaining circle collapse. Their common quotient is the local chart $(\rho_1e^{\mathrm i\theta_1},\rho_2e^{\mathrm i\theta_2})$. Gluing these charts with the prescribed degree-$-1$ normal bundles recovers the degree-five del Pezzo surface $\overline M_{0,5}$. Collapsing all fifteen maximal two-torus fibres is already the first simply connected stage.

For $(1,1)$, the cut space is the compactified modular orbifold with its cusp cut open. The cusp face is $S^1\times B\mu_2\to B\mu_2$: the involution interchanges the two branches of the nodal rational curve and acts trivially on their tensor-product smoothing parameter. Its mapping cylinder fills the cusp while retaining that inertia. The resulting analytic stack is $\overline{\mathcal M}_{1,1}\simeq\mathbb P(4,6)$, including the interior inertia of orders $4$ and $6$ and generic inertia $\mu_2$. The boundary meridian maps to the usual Dehn twist $T\in\mathrm{SL}_2(\mathbb Z)$. The presentation $S^4=1$, $S^2=(ST)^3$ shows directly that setting $T=1$ kills the whole group. Thus the cusp attachment gives the compactification's trivial fundamental group; it retains the stack inertia.

For $(1,2)$, the two divisor normalizations are
\[
 [\overline M_{0,4}/\mathfrak S_2]\quad\hbox{and}\quad
 \overline{\mathcal M}_{1,1}\times\overline M_{0,3}.
\]
In the first, $\mathfrak S_2$ exchanges the node branches and fixes the two marked legs. The two maximal graph types are the loop with a rational tail and the two-edge cycle. Their deployed corner fibres are two-torus bundles over their graph-automorphism stacks; their local groups are respectively the groups $D$ and $C$ of B.10. The loop-tail corner has one retained normal direction for each divisor. The cycle corner has two distinct retained-branch incidences with the irreducible divisor, exchanged by its edge-swap automorphism. The face maps forget the smoothed angular coordinates and then use the sewing maps of B.11. Hence both incidences, and their finite automorphisms, occur in the pushout. The local two-coordinate quotient recovers the node-smoothing chart and the global quotient is $\overline{\mathcal M}_{1,2}$. Killing the two maximal meridian lattices kills every Dehn twist by the extension-to-pants argument; this is again the first simply connected stage.

\begin{result}{D.17}{Proved}{The five depth-three diagrams and the role of the four elementary types}
For $(0,5),(0,6),(1,2),(1,3),(2,0)$, the table below gives every boundary stratum of depth at most three, with its cut-component types and graph automorphisms. On the incidence category of deployed flags, the local group and each arrow are given by the following explicit normal-bundle and sewing construction. Its fundamental group is $\pi_1(L_{g,n})$. Consequently it equals $\Gamma_{g,n}$ in the three complexity-three types $(0,6),(1,3),(2,0)$; in the two complexity-two types its map to $\Gamma_{g,n}$ has kernel $F_{\aleph_0}$.

For every stable type, the action complex of groups on $P_2/\Gamma_{g,n}$ has fundamental group $\Gamma_{g,n}$. Its data use strata of dimension at most two. The connected one-dimensional stage supplies generators, the two-dimensional stage supplies relations, and a stratum of dimension three supplies relations among those relations. The four positive-dimensional single-component elementary types are exactly $(0,4),(1,1),(0,5),(1,2)$.
\end{result}
\begin{proof}
Here is the full geometric specification of the local groups, including their extension classes. For a graph $G$, let
\[
 \mathcal S_G=\left[\prod_v\mathcal M_{g_v,n_v}/\operatorname{Aut}(G)\right].
\]
For each edge $e=\{h,h'\}$ its normal line is
$\mathcal N_e=L_h^\vee\otimes L_{h'}^\vee$, of first Chern class
$-\psi_h-\psi_{h'}$. Let $\mathcal T_G$ be the simultaneous unit-circle bundle of these lines, with the graph group permuting the factors. Define $H_G=\pi_1(\mathcal T_G)$. Pulling to the product of cut-surface Teichm\"uller spaces makes this bundle trivial and proves
\[
 1\longrightarrow\mathbb Z^{E(G)}
 \longrightarrow H_G
 \longrightarrow\pi_1(\mathcal S_G)\longrightarrow1.
\]
The bundle itself specifies the extension, including its graph-automorphism lifts. Equivalently $H_G$ is the stabilizer of the chosen multicurve under sewing; opposite boundary twists of a cut edge are identified with the same global node twist.

A flag colours the edges by the stage at which they are contracted. Its group is the subgroup of $H_G$ preserving that colouring. The two edge injections into adjacent vertex groups are induced by sewing the specified boundary pairs and by retaining the other parametrized normal circles. On the node lattices they send a retained edge basis vector to that same edge basis vector. Graph automorphisms are restricted to those preserving the flag, before sewing. Thus multiple incidences are retained even when their contracted graph types coincide. All graph lists, individual contractions, and edge multiplicities are in the supplied JSON computation; the displayed table groups only isomorphic label-permutation patterns.

For completeness, the twisting elements in an ordinary presentation are specified as follows. Take the barycentric face category of the action complex, choose a representative of each orbit of a chain, and for an arrow $\alpha:A\to B$ choose $g_\alpha\in\Gamma_{g,n}$ carrying the chosen representative of $B$ to its designated face of the chosen representative of $A$. Then
\[
 \psi_\alpha(x)=g_\alpha^{-1}xg_\alpha,\qquad
 f_{\beta,\alpha}=g_\beta^{-1}g_\alpha^{-1}g_{\beta\alpha}
 \in H_C\quad(A\xrightarrow{\alpha}B\xrightarrow{\beta}C).
\]
The equality $\psi_\beta\psi_\alpha=\operatorname{Ad}(f_{\beta,\alpha})\psi_{\beta\alpha}$ follows by multiplying these three words. Multiplying four transports gives the cocycle identity, so this specifies the twisting on every triangle. The presentation has all local groups and arrow letters $u_\alpha$, the relations
\[
 u_\alpha x u_\alpha^{-1}=\psi_\alpha(x),\qquad
 u_\beta u_\alpha=f_{\beta,\alpha}u_{\beta\alpha},
\]
and $u_\alpha=1$ on a chosen maximal tree. This construction is the Borel action groupoid of the actual chain complex: equivariant simplexwise contraction gives its classifying-space comparison, so its homotopy fibre over $B\Gamma_{g,n}$ is that chain complex. For the depth-three link diagram D.2 supplies the one-type equivalence with $L_{g,n}$; D.1 proves exactly the group conclusions stated above.

For $P_2$, D.5 gives simple connectedness, and the same Borel fibre sequence gives $\pi_1=\Gamma_{g,n}$. If $d\leq2$, this assertion includes the empty multicurve, whose vertex is the open stratum. The exact map from the nonempty boundary diagram is induced by adjoining that empty multicurve and all incident chains; its kernel in the two complexity-two examples is the free group in D.1. Thus the low-dimensional and punctured-boundary presentations are related by this concrete map.

At a stratum of dimension one, the positive-complexity cut components have type $(0,4)$ or $(1,1)$, with pants factors. At dimension two they are either a single $(0,5)$ or $(1,2)$ component, or two complexity-one components. This follows by summing $3g_v-3+n_v$ and imposing stability, and proves that precisely the four elementary types furnish the local generator and relation data. D.5's connectedness of $P_1$ and simple connectedness of $P_2$ prove their global roles. At residual dimension three, the upper link is a join of curve complexes with total residual complexity three; its homotopy type is a wedge of two-spheres. Coning that link attaches three-dimensional cells up to homotopy and specifies identities among the two-dimensional relations. The single positive-complexity types are $(0,6),(1,3)$ and, only for the entire closed surface, $(2,0)$; the other possibilities are one complexity-two and one complexity-one component, or three complexity-one components. This lists every stratum type. For the three requested dimension-three moduli spaces, the sole dimension-three stratum is the open stratum itself. The two requested dimension-two spaces have no dimension-three stratum.
\end{proof}

The following enumeration was executed. Vertices are numbered from zero; $00$ denotes a loop and repeated entries denote parallel edges. The labelled number counts the individual strata in that label-permutation orbit. The normal-bundle rule, with these exact edges and cut types, specifies each vertex group and every coloured-flag subgroup.

{\small\paragraph{Type $(0,5)$.}
\begin{longtable}{c|r|l|l|l|r}
$q$ & labelled number & genera; leg counts & edge multiset & cut types & $|\operatorname{Aut}G|$\\\hline
1 & 10 & $(0,0);(2,3)$ & $\{01\}$ & $(0,3),(0,4)$ & 1\\
2 & 15 & $(0,0,0);(1,2,2)$ & $\{01,02\}$ & $(0,3),(0,3),(0,3)$ & 1\\
\end{longtable}
\paragraph{Type $(0,6)$.}
\begin{longtable}{c|r|l|l|l|r}
$q$ & labelled number & genera; leg counts & edge multiset & cut types & $|\operatorname{Aut}G|$\\\hline
1 & 15 & $(0,0);(2,4)$ & $\{01\}$ & $(0,3),(0,5)$ & 1\\
1 & 10 & $(0,0);(3,3)$ & $\{01\}$ & $(0,4),(0,4)$ & 1\\
2 & 60 & $(0,0,0);(1,2,3)$ & $\{01,02\}$ & $(0,3),(0,3),(0,4)$ & 1\\
2 & 45 & $(0,0,0);(2,2,2)$ & $\{01,02\}$ & $(0,3),(0,3),(0,4)$ & 1\\
3 & 15 & $(0,0,0,0);(0,2,2,2)$ & $\{01,02,03\}$ & $(0,3),(0,3),(0,3),(0,3)$ & 1\\
3 & 90 & $(0,0,0,0);(1,1,2,2)$ & $\{01,02,13\}$ & $(0,3),(0,3),(0,3),(0,3)$ & 1\\
\end{longtable}
\paragraph{Type $(1,2)$.}
\begin{longtable}{c|r|l|l|l|r}
$q$ & labelled number & genera; leg counts & edge multiset & cut types & $|\operatorname{Aut}G|$\\\hline
1 & 1 & $(0);(2)$ & $\{00\}$ & $(0,4)$ & 2\\
1 & 1 & $(0,1);(2,0)$ & $\{01\}$ & $(0,3),(1,1)$ & 1\\
2 & 1 & $(0,0);(0,2)$ & $\{00,01\}$ & $(0,3),(0,3)$ & 2\\
2 & 1 & $(0,0);(1,1)$ & $\{01,01\}$ & $(0,3),(0,3)$ & 2\\
\end{longtable}
\paragraph{Type $(1,3)$.}
\begin{longtable}{c|r|l|l|l|r}
$q$ & labelled number & genera; leg counts & edge multiset & cut types & $|\operatorname{Aut}G|$\\\hline
1 & 1 & $(0);(3)$ & $\{00\}$ & $(0,5)$ & 2\\
1 & 3 & $(0,1);(2,1)$ & $\{01\}$ & $(0,3),(1,2)$ & 1\\
1 & 1 & $(0,1);(3,0)$ & $\{01\}$ & $(0,4),(1,1)$ & 1\\
2 & 1 & $(0,0);(0,3)$ & $\{00,01\}$ & $(0,3),(0,4)$ & 2\\
2 & 3 & $(0,0);(1,2)$ & $\{00,01\}$ & $(0,3),(0,4)$ & 2\\
2 & 3 & $(0,0);(1,2)$ & $\{01,01\}$ & $(0,3),(0,4)$ & 2\\
2 & 3 & $(0,0,1);(1,2,0)$ & $\{01,02\}$ & $(0,3),(0,3),(1,1)$ & 1\\
3 & 3 & $(0,0,0);(0,1,2)$ & $\{00,01,12\}$ & $(0,3),(0,3),(0,3)$ & 2\\
3 & 3 & $(0,0,0);(0,1,2)$ & $\{01,01,02\}$ & $(0,3),(0,3),(0,3)$ & 2\\
3 & 1 & $(0,0,0);(1,1,1)$ & $\{01,02,12\}$ & $(0,3),(0,3),(0,3)$ & 1\\
\end{longtable}
\paragraph{Type $(2,0)$.}
\begin{longtable}{c|r|l|l|l|r}
$q$ & labelled number & genera; leg counts & edge multiset & cut types & $|\operatorname{Aut}G|$\\\hline
1 & 1 & $(1);(0)$ & $\{00\}$ & $(1,2)$ & 2\\
1 & 1 & $(1,1);(0,0)$ & $\{01\}$ & $(1,1),(1,1)$ & 2\\
2 & 1 & $(0);(0)$ & $\{00,00\}$ & $(0,4)$ & 8\\
2 & 1 & $(0,1);(0,0)$ & $\{00,01\}$ & $(0,3),(1,1)$ & 2\\
3 & 1 & $(0,0);(0,0)$ & $\{00,01,11\}$ & $(0,3),(0,3)$ & 8\\
3 & 1 & $(0,0);(0,0)$ & $\{01,01,01\}$ & $(0,3),(0,3)$ & 12\\
\end{longtable}
}

\paragraph{Concrete group forms in the first three examples.}
For $(0,5)$, each Petersen vertex $i$ has group
\[
 H_i=\langle f_i,h_{ij_1},h_{ij_2},h_{ij_3}\mid
 [f_i,h_{ij_a}]=1,\quad h_{ij_1}h_{ij_2}h_{ij_3}=f_i\rangle.
\]
This is $\mathbb Z\times F_2$; the relation records normal degree $-1$ and fixes the lifts of the three base peripheral loops. The edge group is $\mathbb Z^2$. Plumbing injects its basis as $(f_i,h_{ij})$ at $i$ and $(h_{ji},f_j)$ at $j$. Non-tree edges carry their arrow letters in the preceding presentation. Abelianizing gives the matrix $-I+A$ of D.4, and hence rank $11$, confirming that this link group is the specified extension of $\Gamma_{0,5}$.

For $(0,6)$, the fifteen $2|4$ divisor groups are the normal circle extensions of $\Gamma_{0,5}$, and the ten $3|3$ divisor groups those of $F_2\times F_2$, with Euler classes $-\psi_h-\psi_{h'}$. The $105$ two-edge groups are the two-circle extensions over a tripod, hence abstractly $\mathbb Z^2\times F_2$; their chosen peripheral lifts remain the prescribed normal-line lifts. The $105$ three-edge groups are $\mathbb Z^3$. All flag maps are the specified sewing maps, and the triangles are exactly the nested edge-contraction flags of the enumerated trees.

For $(1,2)$, B.10--B.11 give the completely presented node-only groups
\[
\begin{aligned}
 N&=\langle a,c,r\mid [a,c]=[a,r]=1,\ r^2=1\rangle,\\
 T&=\langle a,b\mid aba=bab\rangle,\\
 D&=\langle a,h\mid[a,h]=1\rangle,\\
 C&=\langle a,c,q\mid[a,c]=1,\ q^2=1,\ qaq^{-1}=c\rangle.
\end{aligned}
\]
Here $N,T$ are the irreducible and separating divisor groups; $D,C$ are the loop-tail and two-cycle groups. Their node lattices are $\langle a\rangle$, $\langle(ab)^6\rangle$, $\langle a,h^2\rangle$, and $\langle a,c\rangle$. The injections $D\to T$ and $D\to N$ are
\[
 a\mapsto a,\ h\mapsto(ab)^3;\qquad
 a\mapsto a,\ h\mapsto ac^{-1}r.
\]
For a branch of $C\to N$ the flag subgroup is $\langle a,c\rangle$ and its injection is the identity; the other branch is its $q$-conjugate. This restriction is necessary because $q$ exchanges the retained and smoothed edges. The map to the open presentation
$\Gamma_{1,2}=B_4(a,b,c)/\langle(abc)^4\rangle$
is
$h=(ab)^3$, $q=(abc)^2$, $r=(aba)^{-1}(abcaba)(aba)$.
The word checks and external-leg-framed versions are proved and executed in B.9--B.11.

The complete finite word dictionary for the genus-zero $P_2$ action presentations is proved in D.22, using B.18--B.19's exact braid normalization.  D.23 gives the entire fifteen-cell diagram for $(1,2)$, and D.27 gives all nineteen cells and every face word for $(2,0)$.  D.29 completes every cell, face word and twist for $(1,3)$.  The additional compact higher-identity comparisons are recorded explicitly as Open in D.24.  B.20--B.21 separately give the complete sixty-core-path presentation of the bipyramid--pyramid model.


\subsection{The exact graph, boundary, and tropical comparison}

\begin{result}{D.18}{Proved}{The corrected maquette and its comparison fibres}
Let $\mathsf G_{g,n}^{+}$ be the category of stable labelled graphs with at least one edge, with specified contractions and graph isomorphisms as morphisms. Define
\[
 \mathcal B_p=
 \{\,(G,\lambda):|E(G)|=p+1,\ 
       \lambda:E(G)\xrightarrow{\sim}\{0,\ldots,p\}\,\}/\cong .
\]
Relabelling gives the symmetric-group action; an injection of label sets gives a face by contracting precisely the omitted edges. Then $\mathcal B$ is a symmetric semisimplicial set, its realization is the boundary complex, and the map assigning the label weights to the edge lengths is a homeomorphism
\[
 |\mathcal B|\xrightarrow{\sim}\Delta_{g,n}.
\]
Here $\Delta_{g,n}$ is the space of stable tropical curves of total finite-edge length one. The category with the automorphisms retained has the explicit comparison
\[
 B\mathsf G_{g,n}^{+}\simeq
 \operatorname*{hocolim}_{(\mathsf G_{g,n}^{+})^{\rm op}}\Delta_G
 \longrightarrow
 \operatorname*{colim}_{(\mathsf G_{g,n}^{+})^{\rm op}}\Delta_G
 =\Delta_{g,n},\qquad
 \Delta_G=\{(\ell_e)\geq0:\sum\ell_e=1\}.
\]
The actual point fibre of this homotopy-colimit model over a metric graph $(G,\ell)$ has homotopy type $B\operatorname{Aut}(G,\ell)$. The map need not be a quasifibration; this assertion concerns its actual point fibres. Adding the zero-edge graph makes the full graph category contractible, since that graph is terminal.

The analytic boundary end has a compatible comparison to these models. Over a positive-length metric multicurve $(\sigma,\ell)$, the point fibre of its Borel-model map to $\Delta_{g,n}$ is
$B\operatorname{Stab}_{\Gamma_{g,n}}(\sigma,\ell)$.
The map of that fibre to the graph-model fibre is induced by the graph-action homomorphism, whose kernel is the node-framed pure cut-surface group. In particular its homotopy fibre is the classifying space of that kernel.
\end{result}
\begin{proof}
An edge contraction is specified by its contracted edge set and the induced quotient of vertices, half-edges and genera. Consecutive contractions compose by contracting the union of the two sets, transported through the first quotient. Thus injections of finite label sets give contravariant face operations and satisfy their functor identities exactly. Relabelling commutes with these operations. This constructs $\mathcal B$ and proves all its symmetric semisimplicial identities.

A point in its realization has an edge-labelled graph and nonnegative coordinates of sum one, modulo label permutations and the face relations. Contracting exactly the zero-coordinate edges gives a metric stable graph with positive edge lengths. Conversely, choose an edge labelling of a metric graph and use its lengths as simplex coordinates; changing that choice is precisely the symmetric identification. These operations are inverse. The topology on both sides is the final topology from the finitely many graph simplices, so the bijection and its inverse are continuous. In a node-smoothing chart the branches of the boundary are exactly the edges of the dual graph, and simultaneous normalization marks precisely the chosen branches. The boundary-complex face operation is therefore the same edge contraction, including graph monodromy. This gives the asserted identification with the boundary complex. It is also the construction in Chan--Galatius--Payne, Definitions~2.3--2.4 and~3.1 and Theorem~6.1 \cite{D_CGPMarked2022}; we opened the definitions and proof. The unmarked statement and descent proof are Corollaries~5.6--5.7 of \cite{D_CGP2021}.

Every $\Delta_G$ is nonempty and contractible, so the map from its homotopy colimit to the classifying space of the index category is a homotopy equivalence, by the levelwise bar-construction comparison with a constant point diagram. A contraction induces the inclusion of its target simplex as the corresponding zero-coordinate face, and isomorphisms act by their actual coordinate permutations. The ordinary colimit is therefore exactly the quotient just constructed. Projection from the bar construction to that colimit is the displayed map.

Fix a metric graph $x=(G,\ell)$ with all lengths positive. The fibre of that bar-construction map is the nerve of the category of its lifts: objects are graph-simplex representatives mapping to $x$, and arrows are the corresponding face maps and isomorphisms preserving the metric point. Contract all zero-length edges in a lift. This gives a functor to the groupoid of positive-length representatives of $x$, which is equivalent to the one-object groupoid $\operatorname{Aut}(G,\ell)$. The canonical zero-edge contractions give a natural transformation from its inclusion followed by contraction to the identity of the lift category. Its nerve is thus homotopy equivalent to $B\operatorname{Aut}(G,\ell)$. This proves the fibre calculation, including automorphisms of zero-length refinements, which disappear through that natural transformation. For example, at $(1,1)$ the unique nonempty graph is a loop with its flip automorphism: the map is $BC_2\to\{\text{point}\}$. This computes the failure of the original asserted equality of classifying spaces. The zero-edge graph, on the other hand, has a unique contraction from every graph, proving the full category's contractibility.

For the analytic comparison, represent the end as
$E\Gamma_{g,n}\times_{\Gamma_{g,n}}|\mathcal C_{g,n}|$
as in D.1. Send a weighted multicurve to its dual graph with those weights as edge lengths. The classification of marked cut surfaces shows that the orbits are exactly the metric stable graphs, so the orbit space is the same $\Delta_{g,n}$. The preimage of one orbit in the Borel construction is
$E\Gamma_{g,n}\times_{\Gamma_{g,n}}\Gamma_{g,n}/H\simeq BH$,
where $H=\operatorname{Stab}(\sigma,\ell)$.
Cutting the surface and forgetting its marking gives a functor to the metric-graph lift category and hence the compatible comparison above. On this fibre its homomorphism is
\[
 1\longrightarrow K_G\longrightarrow H
 \longrightarrow\operatorname{Aut}(G,\ell)\longrightarrow1,
 \qquad
 1\longrightarrow\mathbb Z^{E(G)}
 \longrightarrow K_G
 \longrightarrow\prod_v\Gamma_{g_v,n_v}\longrightarrow1.
\]
Surjectivity onto graph automorphisms follows by choosing oriented homeomorphisms between the cut components with the prescribed flag labels and sewing them. The kernel consists exactly of mapping classes fixing each component and each half-edge; the cut-surface sewing sequence gives its second displayed extension. Thus $BK_G$ is the homotopy fibre of the stated map $BH\to B\operatorname{Aut}(G,\ell)$. The normal Euler classes are those in D.17 and A.12.
\end{proof}

The terminology in Lead~6 can now be assigned exact maps. Specified MD-graph generizations, with their graph automorphisms, give $\mathsf G_{g,n}$ and its nerve. Taking isomorphism classes of edge-labelled graphs gives $\mathcal B$. Taking isomorphism classes of strictly contracting flags gives the barycentric subdivision of $\mathcal B$: on each graph simplex this is the ordinary subdivision by nested nonempty coordinate faces, and the construction commutes with every graph isomorphism and contraction. Its realization maps homeomorphically to $|\mathcal B|$ by sending a face-chain vertex to the corresponding simplex barycenter. This describes the comparison without a degreewise identification of edge-labelled cells with contraction chains. The levelwise component sets of the original groupoid-valued flag data are addressed in A.11 and A.13; their lost middle-endpoint identifications need not satisfy a Segal composition law.

For the two examples in the added task, this map has concrete consequences. For the tripod, the tropical boundary is three points; the analytic end's three circle fibres are its three meridian groups, while the trivially valued tube adds a positive radial interval at each point. For $\overline M_{0,5}$, graph automorphisms are trivial and the graph model is the Petersen graph. The analytic end maps to it with circle-bundle groups over vertex interiors and $\mathbb Z^2$ over edge interiors; the nonarchimedean tube retracts to it by D.4. The exact maps, the Betti numbers in D.4, and the proper-modification comparisons in D.3 establish all asserted relationships in these examples.


\subsection{The full low-genus completion range and two source corrections}

\begin{result}{D.19}{Known; transfer proved}{Congruence comparison in genus at most two}
For every stable type of genus $0,1,$ or $2$, the natural map
$\widehat\Gamma_{g,n}\to\check\Gamma_{g,n}$
is an isomorphism. Consequently the full finite-level and procongruence completions of every $P_k$ are canonically isomorphic in these genera. In particular, their $P_1$ are connected and their $P_k$, $k\geq2$, have trivial profinite fundamental group. The full unmodified curve complex is connected with trivial profinite fundamental group in this range exactly when $d\geq3$; the two connected complexity-two exceptions have nontrivial profinite fundamental group by D.12 and D.21.
\end{result}
\begin{proof}
The genus-zero congruence theorem is the puncture induction in Boggi's Proposition~2.6, reproducing Asada's argument: start with $\Gamma_{0,3}=1$, use the completed Birman sequence, and use the tangential lift to identify the point-pushing closure with the full profinite surface group \cite{D_Boggi2009}. The opened Theorem~3.5 and Lemmas~3.6--3.7 of the same paper prove faithfulness for the full hyperelliptic group with at least one puncture. In genera one and two the hyperelliptic locus is the entire moduli stack, so this gives every stable genus-one type and every genus-two type with $n\geq1$. The closed genus-two case requires an additional argument; it is not part of that open-surface theorem.

Here is the independent closed-surface transfer, which avoids both problematic claims discussed in D.20. Write
\[
 P=\widehat{\pi_1(S_2)},\qquad
 Q=\widehat{\pi_1(\operatorname{Conf}_2(S_2))},\qquad
 F=\widehat{\pi_1(S_2\setminus\{x\})}.
\]
The full universal-curve and configuration sequences give
\[
 1\longrightarrow Q\longrightarrow\widehat\Gamma_{2,2}
       \longrightarrow\widehat\Gamma_{2,0}\longrightarrow1,
 \qquad
 1\longrightarrow F\longrightarrow Q\longrightarrow P\longrightarrow1.
\]
These are sequences of full geometric fundamental groups, obtained from the ordered configuration family, as in Part~E; no congruence exactness is assumed. Its geometric monodromy preserves the forgetful fibrations and diagonal inertia, so its image in $\operatorname{Out}(Q)$ is FC-admissible in the sense of Hoshi--Mochizuki. The injectivity assertion of their Theorem~6.1, with all primes and $n=1$, applies to this closed genus-two curve and gives
\[
 \operatorname{Out}^{\rm FC}(Q)\hookrightarrow
 \operatorname{Out}^{\rm FC}(P).
\]
We opened that theorem and its proof, including its reduction to Corollary~5.3; it allows a proper hyperbolic curve at this value of $n$ \cite{D_HoshiMochizuki2011}.

Let $y\in\widehat\Gamma_{2,0}$ act trivially in $\operatorname{Out}(P)$. The displayed injection makes its outer action on $Q$ trivial. Lift $y$ to $z\in\widehat\Gamma_{2,2}$. Conjugation by $z$ on $Q$ is inner, so multiplication by a suitable element of $Q$ makes $z$ centralize $Q$, without changing its image $y$. In particular $z$ centralizes the distinguished subgroup $F$. Its image $z_1\in\widehat\Gamma_{2,1}$ therefore acts trivially in $\operatorname{Out}(F)$: this action is precisely restriction of conjugation by a lift in $\widehat\Gamma_{2,2}$. The already established open genus-two congruence theorem forces $z_1=1$, and its image $y$ is also $1$. This proves the closed case, with the exact groups and the necessary conjugator adjustment specified.

An isomorphism of these two completions means exactly that finite-index congruence subgroups are cofinal among all finite-index subgroups. Restricting the same quotient-complex diagram to a cofinal index system gives the asserted isomorphism of profinite complexes. D.11 then gives the stated connectedness and fundamental groups. At complexity one the unmodified curve complex is disconnected; D.12 and D.21 settle complexity two. This argument asserts no unproved higher profinite connectivity.
\end{proof}

\begin{result}{D.20}{Refuted}{Two literal centralizer assertions in the retrieved Boggi source}
Lemma~6.7 of the retrieved 12 July 2021 version of \emph{Congruence topologies on the mapping class group} asserts that the hyperelliptic involution is self-centralizing in every finite covering-transformation semidirect product specified there. The characteristic mod-two homology cover refutes that assertion. Corollary~4.3 of the same version, when stated for arbitrary unequal multitwist exponents as printed, also cannot identify the centralizer with the full stabilizer of the underlying curve set. Neither assertion is used in D.19.
\end{result}
\begin{proof}
For a closed hyperelliptic surface of genus $g\geq2$, choose a Weierstrass basepoint and set
\[
 N=\ker\bigl(\pi_1(S_g)\longrightarrow H_1(S_g;\mathbb F_2)\bigr).
\]
This subgroup is proper, normal, finite index and characteristic, so it satisfies every subgroup hypothesis in the printed Lemma~6.7 \cite[pp.~26--27]{D_Boggi2020}. The hyperelliptic involution acts as $-I$ on integral homology and hence trivially on the quotient
$G_N=(C_2)^{2g}$. Thus the specified deck group is
\[
 G=G_N\rtimes\langle\iota\rangle
   =(C_2)^{2g+1},\qquad C_G(\iota)=G.
\]
Its centralizer has order $2^{2g+1}$, whereas $\langle\iota\rangle$ has order two. The unramified surface cover has degree $2^{2g}$ and genus $1+2^{2g}(g-1)$ by Euler characteristic. The saved exact computation gives degree $16$, genus $17$, and centralizer order $32$ at $g=2$.

For the second assertion choose disjoint nonseparating curves $a_1,a_2$ on the two handles of a closed genus-two surface, and an orientation-preserving mapping class $s$ exchanging the handles. Then $s$ belongs to the setwise stabilizer of $\{a_1,a_2\}$, which is the stabilizer explicitly defined in Theorem~3.6 of that source. But
\[
 s(T_{a_1}T_{a_2}^{\,2})s^{-1}
     =T_{a_1}^{\,2}T_{a_2}
     \ne T_{a_1}T_{a_2}^{\,2}.
\]
The inequality is visible already on $H_1(S_2;\mathbb F_3)$: in the symplectic basis $(a_1,b_1,a_2,b_2)$ the two matrices have upper-right entries $(1,2)$ and $(2,1)$ in their two diagonal $2$-by-$2$ blocks. Thus it persists in the procongruence group, without any appeal to a centralizer theorem. Both examples were executed in \path{code/d_hyperelliptic_counterexample.py}. The first refutes a universal finite-quotient statement; it does not refute the separate profinite self-normalizer assertion, which requires a cofinal argument. The second shows exactly why a weight-preserving stabilizer must appear when unequal exponents are used.
\end{proof}

\begin{result}{D.21}{Proved}{The second complexity-two profinite obstruction}
The full finite-level completion of the unmodified curve complex $\mathcal C_{1,2}$ has nonzero continuous first cohomology with $\mathbb F_2$ coefficients, and nontrivial profinite fundamental group. By D.19 the same holds for its procongruence completion.
\end{result}
\begin{proof}
Use the torus involution exchanging the two marked points. Its quotient is a sphere with four branch points and a fifth distinguished point, the image of the two marks. Lifting a sphere curve gives an isomorphism
\[
 P^*:\mathcal C_{0,5}\xrightarrow{\sim}\mathcal C_{1,2}.
\]
The precise map and invariant-representative proof are Luo, Lemma~2.1(a), opened on pages~3--4 \cite{D_Luo2000}. A sphere curve separating two branch points lifts to two parallel nonseparating curves and hence one isotopy class. A curve separating the distinguished point and one branch point lifts to one separating curve bounding the disk containing both torus marks. Hyperelliptic invariance supplies an invariant representative of every torus curve; taking invariant disjoint representatives proves the converse disjointness implication. These operations give the displayed inverse correspondence.

The construction is equivariant for the lift of the pure sphere mapping class group. Algebraically, B.9 identifies
\[
 \Gamma_{1,2}=B_4/\langle\Delta_4^2\rangle,\qquad
 \Gamma_{0,5}=P_4/\langle\Delta_4^2\rangle,
\]
and the latter subgroup has index $24$, by the permutation homomorphism to $\mathfrak S_4$. The subgroup is exactly the one retaining all four branch labels in the preceding cover construction. Finite-index subgroups of this subgroup are cofinal in the finite-index system for $\Gamma_{1,2}$; normality can be imposed by taking the finite intersection of conjugates. The equivariant simplicial isomorphism therefore identifies the two full finite-level curve complexes. Transporting the explicit continuous Petersen-edge class and the closed pentagon of D.12 gives a nonzero continuous $\mathbb F_2$ class on the genus-one complex. Its profinite fundamental group is consequently nontrivial. D.19 supplies the final procongruence comparison.
\end{proof}


\subsection{The named finite complexes on the second level}

\begin{result}{D.22}{Proved}{Complete word data on $P_2$ in the five- and six-point sphere cases}
The complexes of groups on $P_2/\Gamma_{0,n}$ for $n=5,6$
have the following complete finite specifications:
\[
\begin{array}{c|rrrrrr}
n&\text{vertices}&\text{arrows}&\text{triangles}&
 \text{local generators}&\text{generator injections}&
 \text{tree arrows}\\\hline
5&26&55&30&91&175&25\\
6&235&840&630&1110&3570&234.
\end{array}
\]
Every vertex presentation, every image of an edge-group generator,
every transporter and every triangle twisting element is an exact
integer word in the certificate
\path{code/d_p2_genus_zero.json}.  The definitions below give the
same data without a choice left unspecified.  Their fundamental
groups are $\Gamma_{0,5}$ and $\Gamma_{0,6}$, respectively.

The maps to the presentations P5c2 and P5c3 are given on every
listed generator by that certificate.  Their inverses send
\[
\begin{array}{c|l}
5&a,b,x,y,z\longmapsto
 A_{13},A_{23},A_{14},A_{24},A_{34},\\
6&p,q,r,s\longmapsto A_{15},A_{25},A_{35},A_{45},
\end{array}
\]
where each $A_{ij}$ on the right is the specified local node
generator at the divisor with cluster $\{i,j\}$, and the second
row supplements the first.  Thus the comparison is proved
generator by generator.  There is no dimension-three stratum
in $M_{0,5}$.  The unique such stratum for $M_{0,6}$ is its open
stratum; the certificate gives its $630$ complete-flag twisting
coherence identities.
\end{result}
\begin{proof}
We specify both the finite index set and all words.  The labelled
graphs and their indices are those of the exhaustive files
\path{code/a_graphs_0_5.json} and
\path{code/a_graphs_0_6.json}.  For five labels keep every tree,
including the zero-edge tree.  For six labels keep precisely
the nonempty trees.  There is one arrow $G\to H$ for each
strict contraction, and one triangle for each strict chain of
two contractions.  A labelled stable genus-zero tree has no
nontrivial graph automorphism: an automorphism fixing its leaves
fixes every internal vertex as the intersection of three leaf
paths.  Its edges are distinguished by their labelled splits.
Consequently these trees and chains are exactly the respective
orbits of vertices, edges and triangles of $P_2$, with no
unrecorded parallel incidence.

Put $m=n-1$ and root every tree at the vertex incident to the
last label, denoted $\infty$.  The other labels are
$1,\ldots,m$.  At each vertex order the children by the least
label in their descendant leaves.  Reading the leaves in this
depth-first order gives a permutation $\pi_G$.  Every cluster
away from the root is an interval in this order.  Let $P_G$
be the positive permutation braid found by the following fixed
algorithm: read $\pi_G$ from left to right, and move its next
entry left to its prescribed position by adjacent positive
crossings.  In that ordered configuration take round circles
around all the clusters, and pull them back by $P_G^{-1}$.
This is the chosen multicurve $\sigma_G$.

The proof and certificate use braid words in traversal order.
Write
\[
 A_{ij}=\sigma_{j-1}\cdots\sigma_{i+1}\sigma_i^2
          \sigma_{i+1}^{-1}\cdots\sigma_{j-1}^{-1},
 \qquad Z_q=(\sigma_1\cdots\sigma_{q-1})^q.
\]
At a vertex $v$ of $G$, let $q_v$ be its number of children,
where a marked leaf is one child and a descendant subtree is
one child.  Thus $q_v=\operatorname{val}(v)-1\geq2$.
The round-system stabilizer in $PB_m$ is
\[
                  \prod_{v\in V(G)}PB_{q_v}.                 \tag{D.1}
\]
Here is its actual embedding.  Replace every strand at $v$ by
the block of its descendant leaves, retaining their fixed
order.  Cable every crossing with zero rotation of those
blocks, and extend it by the identity outside the interval of
the vertex.  If adjacent blocks have sizes $a,b$, their positive
exchange is the word obtained by moving the $b$ strands of the
right block, successively from left to right, across all $a$
strands of the left block.  Its explicit shifted word is
\[
 \prod_{j=0}^{b-1}
       (\sigma_{s+a+j}\sigma_{s+a+j-1}\cdots\sigma_{s+j+1}),
                                                               \tag{D.2}
\]
where $s$ is the number of strands to its left.  The negative
exchange is the inverse positive exchange with block sizes
$b,a$.  Update the ordered blocks after each crossing.
This defines the cable of every signed Artin word.

To justify (D.1), collapse the outermost child disks to points.
A pure stabilizer gives a pure braid of these points and an
internal braid in each disk.  The motion of disks with their
coordinates transported by translation gives a section: choose
their radii smaller than one third of the minimum separation
of their centres.  This keeps the coordinates unrotated.
After removing that motion, the remaining homeomorphisms are
supported inside the child disks.  They commute with that
zero-rotation motion up to isotopy, because the coordinates in
each transported disk identify their supports throughout the
isotopy.  Thus the group is the direct product of the outer
pure braid group and its internal stabilizers.  Induction on
the rooted tree proves (D.1), including injectivity and
surjectivity of the specified cable map.

The full turn of all $m$ strands is the product of the local
full turns $Z_{q_v}$.  Indeed the outer full turn winds every
pair of distinct child disks once around one another and the
internal full turns supply exactly the turns within the
children; the disk motions give the asserted braid equality.
The program independently checks it by the faithful Artin
action for all $261$ trees.  Passing from the disk braid group
to the sphere mapping class group quotients by this full turn.
The actual local group is therefore
\[
 H_G=\left(\prod_vPB_{q_v}\right)
        \big/\left\langle\prod_v Z_{q_v}\right\rangle.        \tag{D.3}
\]
Its generator at $v$ is named $(v,i,j)$ and embedded in the
global group by the word
$P_G\,\operatorname{cable}_v(A_{ij})\,P_G^{-1}$.

For a complete finite presentation of each factor, use
$PB_q=F_{q-1}\rtimes PB_{q-1}$ successively.  For each
$i<j<\ell\leq q$ and each $k<\ell$, impose
\[
 A_{ij}A_{k\ell}A_{ij}^{-1}
       =\operatorname{comb}_{\ell}
                  (A_{ij}A_{k\ell}A_{ij}^{-1}).               \tag{D.4}
\]
On the right use the exact pure-braid combing of B.18; the
word lies in the free generators $A_{1\ell},\ldots,
A_{\ell-1,\ell}$.  The punctured-plane fibration and its
section prove this presentation inductively, by the same
normal-form argument as B.19.  Adjoin commutators between
distinct factors, and the single diagonal central relation
in (D.3).  The certificate expands all these words, so (D.4)
is an algorithm specifying actual finite relations and not a
reference to unknown conjugations.

Now take a contraction $G\to H$.  Every $H$-cluster is an
interval in both $\pi_H$ and $\pi_G$.  Apply the same positive
permutation algorithm to change the former ordered leaves
into the latter.  Denote its word by $B_{HG}$.  It transports
the round $H$-system to the corresponding subsystem of the
round $G$-system.  To see this directly, any strand outside
one of the blocks either crosses every strand of that block
once in the same direction or crosses none; the block can
therefore be transported as a disk.  The same statement holds
recursively for its subblocks.  The pure transporter, in the
fixed traversal convention, is
\[
                q_{GH}=P_H B_{HG}P_G^{-1}.                   \tag{D.5}
\]
It carries $\sigma_H$ to the retained part of $\sigma_G$.
Whenever those two multicurves already agree, choose instead
the empty word.  This last decision is exact: compute their
boundary words in the free punctured-disk group, cyclically
reduce them and compare up to inversion.  Free homotopy of
essential simple curves is isotopy.  The computation checks
every claimed disk transport on those same boundary words.

The edge injection is $x\mapsto q_{GH}xq_{GH}^{-1}$.
It remains to express that answer in the named target factors,
rather than merely in the ambient group.  In the round
$H$-configuration, delete every strand except the first one
in each child block.  This reads the outer pure braid exactly.
Remove its zero-rotation cable, restrict the residual braid to
each child disk, and repeat.  This is inverse to (D.1), so it
terminates and gives the target word in its local $(v,i,j)$.
The certificate gives all $175$ and $3570$ generator images
obtained this way, and checks each reconstructed braid using
the faithful Artin action.

For $G\to H\to J$ put
\[
       f_{GHJ}=q_{HJ}q_{GH}q_{GJ}^{-1}\in H_J.               \tag{D.6}
\]
It fixes the representative $\sigma_J$ because the first two
transports and the direct transport have the same retained
multicurve.  Apply the preceding inverse cable algorithm to
give its exact word in $H_J$.  The local group conjugations
and (D.6) prove the twisting identity.  For a further
contraction to $K$ they also give
\[
 f_{HJK}f_{GHK}
       =\psi_{JK}(f_{GHJ})f_{GJK},                           \tag{D.7}
\]
by multiplying the actual transport words.

There are respectively $54$ and $789$ arrows whose chosen
transporter is the identity.  The executed breadth-first
search proves that these arrows connect the entire finite
vertex set.  The certificate records a maximal tree using
only these arrows.  Therefore the standard complex-of-groups
presentation uses every local presentation above, one letter
$u_{GH}$ for each arrow, and the explicit relations
\[
 u_{GH}x u_{GH}^{-1}=\psi_{GH}(x),\qquad
 u_{HJ}u_{GH}=f_{GHJ}u_{GJ},\qquad
 u_{GH}=1\quad\text{on the recorded tree}.                   \tag{D.8}
\]
The map to the ambient mapping class group sends local
generators to their cable words and $u_{GH}$ to $q_{GH}$;
the tree relations hold literally.  These are the genuine
stabilizers and transports of the action on $P_2$.
The Borel fibre sequence and D.5 therefore prove that this
map on fundamental groups is an isomorphism.

Apply B.18's exact normalization to every ambient braid word
in (D.8), discard the exponent of the global full turn, and
use B.19's named generators.  This is the forward dictionary
recorded for every local generator, arrow and twist.  For
the inverse, the chosen representative of the pair divisor
$\{i,j\}$ has node twist exactly $A_{ij}$: its permutation
braid is $\sigma_{j-1}\cdots\sigma_{i+1}$ on that block,
and conjugating $\sigma_i^2$ gives the displayed band word.
The program verifies these five and nine identities and
records the precise graph and local-generator indices.
The six and twenty-six relations of B.19 hold on these
inverse images because the independently proved ambient
map is injective.  The composites are the identity on every
named generator, again by that injection.  This proves both
directions of the required presentation comparison.

Finally, in the six-point case adjoin the empty multicurve,
the sole dimension-three stratum.  All its transporters are
the identity.  For every complete flag $G\to H\to J\to0$
the identity (D.7) is consequently the specified
tetrahedral coherence among the triangle relations; the
file lists all $630$ flags and both sides in the compact
moduli generators.  These are the identities supplied by
the open stratum in this finite action diagram.  The
corresponding unquotiented upper link is the curve complex
$\mathcal C_{0,6}$, a wedge of two-spheres, so its cone
attaches precisely the dimension-three homotopy data
described in D.17.
\end{proof}

\begin{result}{D.23}{Proved}{The complete fifteen-cell second-level diagram for $(1,2)$}
Let $O,N,T,D,C$ be the empty, nonseparating, separating,
loop-tail and two-cycle multicurve orbits, respectively.
The quotient of $P_2$ has five vertex orbits, seven edge
orbits and three triangle orbits.  Its barycentric face
complex of groups has exactly the following fifteen cells:
\[
\begin{array}{c|l}
\text{cell group}&\text{cells}\\\hline
\Gamma_{1,2}&O\\
N&N,\ NO\\
T&T,\ TO\\
D&D,\ DN,\ DT,\ DO,\ DNO,\ DTO\\
C&C,\ CO\\
A&CN,\ CNO.
\end{array}
\]
The groups $N,T,D,C$ and their exact generators are those
displayed in D.17, and
$A=\langle a,c\mid[a,c]=1\rangle$.
Every face transporter and every triangle twist is the
identity.  Every face injection is specified below and
listed individually in
\path{code/d_p2_elliptic_two.json}.  There are $32$ face
arrows and $18$ face triangles.  The resulting fundamental
group has precisely the presentation
\[
 \langle a,b,c\mid [a,c]=1,\ aba=bab,\ bcb=cbc,
                         (abc)^4=1\rangle.
\]
The three generators come from the one-dimensional strata:
$a,c$ from $N$, and $b$ from $T$.  The open stratum,
of dimension two, supplies the displayed global relations.
There is no dimension-three stratum.
\end{result}
\begin{proof}
Use the actual curve representatives
\[
 O=\varnothing,\quad N=\{a\},\quad T=\{d\},\quad
 D=\{a,d\},\quad C=\{a,c\},\qquad d=(ab)^6.
\]
Here a curve is named by its positive twist.  B.9--B.11
prove the geometric stabilizer identifications, their
presentations and their embeddings into
$\Gamma_{1,2}=B_4/\langle\Delta_4^2\rangle$.
Every flag in the displayed list is a literally nested
flag of these representatives.  Thus all face transports
can be the identity.  The two possible retained members
of $C$ give one orbit, interchanged by $q$; its chosen
flag has stabilizer $A$, a proper subgroup of index two
in $C$.  In particular $CN\to C$ and $CNO\to CO$ use
$A\hookrightarrow C$.  This retains both incidences in
the action diagram.  All other edge and triangle
stabilizers are the displayed full intersections.

Here are all nonidentity types of subgroup injections;
maps between cells with the same group are identities.
The maps $D\to T$ and $D\to N$ are
\[
 (a,h)\longmapsto(a,(ab)^3),\qquad
 (a,h)\longmapsto(a,ac^{-1}r).
\]
The maps $A\to N$ and $A\to C$ retain $a,c$.
All maps to $O$ are given by
\[
 a\mapsto a,\quad b\mapsto b,\quad c\mapsto c,
 \quad h\mapsto(ab)^3,\quad q\mapsto(abc)^2,
 \quad r\mapsto(aba)^{-1}(abcaba)(aba).
\]
Their composites agree: the only additional named equality
is $(ab)^3=ac^{-1}r$, proved as an exact braid word in
B.11.  Thus these prescriptions specify all $32$ face
arrows and all $18$ commuting face triangles.
The program independently verifies every local relator,
every generator image and every composite using the
faithful Artin action, retaining the central full twist
before quotienting.  Its executed output accompanies the
complete list.

The empty multicurve is a fixed cone point of $P_2$.
Adding $O$ to any boundary flag leaves its stabilizer
unchanged.  In the face diagram, the inclusion to that
cone flag is consequently an isomorphism on its local
group, while the other face map is its displayed
inclusion into $\Gamma_{1,2}$.  Contracting these cone
pairs gives the one-object groupoid of $\Gamma_{1,2}$.
This proves the exact fundamental-group comparison and
its inverse, rather than merely comparing group orders.

For the requested generator assertion, identify the
open letters $a,c$ with their images from $N$, and
$b$ with its image from $T$.  Eliminate $h,q,r$ by the
displayed words.  The cone contraction shows that all
remaining local and face relations are precisely the
relations already holding in the open presentation.
Conversely those four open relations give every local
and face relation, by the checked embeddings.  The two
maps are the identity on $a,b,c$, proving the stated
generator-by-generator equivalence.  The maximum
modular dimension is two, proving the final assertion.
\end{proof}

\begin{result}{D.24}{Open}{The remaining compact higher-identity dictionaries in Added task 4}
D.22 supplies every finite word for $(0,5)$ and $(0,6)$
and proves their group-presentation comparisons with
P5c2 and P5c3.  D.23 gives all named word data for
$(1,2)$, D.27 does so for $(2,0)$, and D.29 completes
$(1,3)$.  Thus all five requested vertex, edge,
transport, sewing and triangle-twist dictionaries
are now supplied, with explicit inverse maps on the
compact global generators.

D.30 now supplies a defined cellular lift of the
six-point comparison: all local and sewing relator
fillings, all $630$ geometric tetrahedral identity
sequences, their Peiffer reductions and coordinates,
and an inverse inclusion of the $24$ open prism
cells.  Its selected degree-three matrix is
$[0\mid I_{24}]$ in D.25's orientation.  These are
coordinates for the specified mapping-cylinder
model, whose entire lower-cell comparison is given.

The remaining compact higher-identity task concerns
positive genus.  D.27 and D.29 give, respectively, all
$24$ genus-two and $198$ three-pointed-torus cone
coherences.  Their expansions in the free crossed
modules of the corresponding compact global
presentations remain additional higher-identity
comparisons.  No missing finite stabilizer or
sewing-word datum is being left implicit in this
remaining scope.
\end{result}

\begin{result}{D.25}{Proved}{The twenty-four compact identity words and their basis map}
Let $F$ be the free group on the nine generators
$a,b,x,y,z,p,q,r,s$, and let $\mathcal R$ be the ordered
list of the six base relators and twenty conjugation
relators of B.19.  The program
\path{code/d_p5c3_identities.py} expands all twenty-four
identities indexed by a base relator $R$ and one of
$u=p,q,r,s$ as products
\[
                  \prod_\nu c_\nu r_\nu^{e_\nu}c_\nu^{-1}=1
                  \quad\text{in }F,\qquad
                  r_\nu\in\mathcal R,\ e_\nu=\pm1.
\]
Every conjugator, relator index and sign is listed in
\path{code/d_p5c3_identities.json}.  There are $2218$
conjugate factors in all, and the largest identity has
$640$ factors.  Every expanded product was freely reduced
to the empty word.

Under the standard map from identity sequences modulo
Peiffer identities to the second homotopy group of the
presentation complex, these twenty-four sequences form
a free $\mathbb Z[\Gamma_{0,6}]$-basis.  Their exact
comparison with the twenty-four sequences of P5c3 is
inversion, separately for each pair $(R,u)$; on this
module the comparison is multiplication by $-1$ on
each named basis vector.
\end{result}
\begin{proof}
Write $\phi_w$ for the actual free-group action in B.19,
including its computed inverse automorphisms, and put
\[
                   D_w(v)=wvw^{-1}\phi_w(v)^{-1}.
\]
For positive base and fibre generators $h,u$, this is
one of the twenty mixed relators.  All other cases
are expanded by the following exact free-word identities:
\[
\begin{split}
D_{ab}(v)&=aD_b(v)a^{-1}D_a(\phi_b(v)),\\
D_h(vv')&=D_h(v)\phi_h(v)D_h(v')\phi_h(v)^{-1},\\
D_h(v^{-1})&=\phi_h(v)^{-1}D_h(v)^{-1}\phi_h(v),\\
D_{h^{-1}}(v)&=h^{-1}D_h(\phi_h^{-1}(v))^{-1}h.
\end{split}                                                   \tag{D.9}
\]
Each equation follows by cancellation of the displayed
words; it fixes all conjugators and signs.  The program
uses the first rule recursively along the base word,
and the other three rules along the fibre word.
For each of the six base relators, B.19's checked action
has $\phi_R(u)=u$ freely.  Hence the identity sequence
constructed here is
\[
              D_R(u)\,(uRu^{-1})\,R^{-1}=1.                  \tag{D.10}
\]
It is exactly the inverse sequence to P5c3's
$R(uR^{-1}u^{-1})D_R(u)^{-1}$, without a change of
generators, relators or base paths.  The final check
expands each factor back into the nine free letters;
it does not use the desired quotient equality to
certify the identity.

Here is the completeness argument, independently
checked against the full attached P5c3 proof.  The
six-relator presentation of $\Gamma_{0,5}$ is the
two-dimensional graph bundle over the two-petalled
rose, with fibre a three-petalled rose.  Its universal
cover is a tree of contractible fibre trees, so its
presentation complex $K_5$, with cell counts $(1,5,6)$,
is aspherical.  The exact action on the free group
$F(p,q,r,s)$ builds a graph bundle up to homotopy over
$K_5$ with fibre the four-petalled rose $R_4$ and a
section.  To make this construction explicit, use
the four-edge graph maps over every base edge; over
each base two-cell use the free reductions
$\phi_R(u)=u$ for its four edges.  There is no further
obstruction or unspecified homotopy class: each
component of
\[
                 \operatorname{Map}_*(R_4,R_4)
                         \simeq(\Omega R_4)^4
\]
is contractible.  It supplies the compatible homotopies
for these maps.  Equivalently, gluing their mapping
cylinders over the cells gives the homotopy bundle
classified by that precise action.

Over a base $q$-cell attach the section cell first.
The remaining relative cell pair is
\[
 (D^q\times R_4,\ S^{q-1}\times R_4\cup D^q\times\{*\}),
\]
whose quotient is $S^q\wedge R_4$, a wedge of four
$(q+1)$-spheres.  Thus the total complex has cell counts
$(1,9,26,24)$ and no higher cells.  Its two-skeleton
is exactly B.19's presentation complex.  The boundary
of the prism over $(R,u)$ is (D.10), with its specified
orientation.  The homotopy bundle has aspherical base
and fibre, so its homotopy exact sequence gives
asphericity and the group $\Gamma_{0,6}$.

Passing to the universal cover of this three-complex
therefore gives the exact sequence
\[
 0\longrightarrow\mathbb Z[\Gamma_{0,6}]^{24}
  \xrightarrow{\partial_3}\mathbb Z[\Gamma_{0,6}]^{26}
  \xrightarrow{\partial_2}\mathbb Z[\Gamma_{0,6}]^9.
\]
The universal cover of its two-skeleton is simply
connected.  Hurewicz identifies its $\pi_2$ with
$\ker\partial_2$, and $\partial_3$ is an isomorphism
from the indicated free module onto that kernel.
Its columns are precisely the twenty-four attaching
spheres (D.10).  This proves the claimed basis and
its sign comparison.  D.30 supplies the full two-cell
comparison and the coordinates of all $630$ geometric
cone tetrahedra for its specified mapping-cylinder
model, with the inverse inclusion of these prism cells.
\end{proof}


\subsection{A finite symplectic counterexample in the source chain}

\begin{result}{D.26}{Refuted}{The integral conjugacy assertion in Boggi's Lemma 5.12}
Lemma~5.12 on page~26 of the retrieved version~5 of
\emph{On the procongruence completion of the Teichm\"uller modular group}
asserts that two nonconjugate finite cyclic subgroups of
$\operatorname{Sp}_{2g}(\mathbb Z)$ remain nonconjugate modulo
every integer $m\geq3$.  This assertion already fails for
two order-two subgroups of $\operatorname{Sp}_4(\mathbb Z)$
at $m=3$.  The exact counterexample below therefore also
prevents using the printed proof of Lemma~5.10, and its
subsequent use in Section~5, without an additional repair.
This statement does not assert a counterexample to the
separate profinite twist-normalizer conclusions.
\end{result}
\begin{proof}
Put
\[
 J=\begin{pmatrix}0&1\\-1&0\end{pmatrix},\qquad
 \Omega=\operatorname{diag}(J,J),\qquad
 A=\begin{pmatrix}I_2&0\\0&-I_2\end{pmatrix},\qquad
 B=\begin{pmatrix}0&I_2\\I_2&0\end{pmatrix}.
\]
Both $A$ and $B$ have order two and preserve the integral
symplectic form $\Omega$.  For an integral involution $T$
on $L=\mathbb Z^4$, the index
\[
       [L:L^+(T)+L^-(T)],\qquad
       L^\pm(T)=\ker(T\mp I:L\to L),
\]
is invariant under integral conjugacy: a conjugating
integral automorphism carries the two displayed kernels
onto the corresponding kernels for its conjugate.
For $A$ the two coordinate planes give index $1$.
For $B$ the two eigenlattices are
\[
 L^+(B)=\{(v,v):v\in\mathbb Z^2\},\qquad
 L^-(B)=\{(v,-v):v\in\mathbb Z^2\}.
\]
Their sum consists of the pairs $(x,y)$ with
$x\equiv y\pmod2$, and hence has index $4$.
Thus $A$ and $B$ are not conjugate even in
$\operatorname{GL}_4(\mathbb Z)$.  Their order-two cyclic
subgroups are not conjugate either, since their unique
nonidentity elements would have to be conjugate.

Now take
\[
 Q=\begin{pmatrix}
 1&0&1&0\\
 0&2&0&2\\
 1&0&-1&0\\
 0&2&0&-2
 \end{pmatrix}.
\]
Direct multiplication gives
\[
       BQ=QA,\qquad \det Q=16,\qquad
       Q^{\mathsf t}\Omega Q=4\Omega.
\]
Consequently the reduction $\overline Q$ is symplectic
and invertible modulo $3$, and
$\overline Q\,\overline A\,\overline Q^{-1}=\overline B$.
This is the required conjugacy in
$\operatorname{Sp}_4(\mathbb F_3)$.
All matrix products, determinants and eigenlattice
indices were independently executed and saved in
\path{code/d_symplectic_involutions.py}, with its JSON
certificate and text output.

The exact printed lemma and the paragraph immediately
preceding it were opened and read in the source
\cite[p.~26]{ET_Boggi2014}.  That paragraph explicitly
invokes this lemma to finish Lemma~5.10.  The fixed-locus
argument printed for Lemma~5.12 therefore cannot establish
its claimed implication; the two matrices above give
its precise failed conclusion.  The cutting theorem
used in D.19 and Part~E is Theorem~4.5, whose proof
precedes this argument and was checked separately.
\end{proof}


\subsection{The complete genus-two second-level word diagram}

\begin{result}{D.27}{Proved}{Nineteen cells and every sewing word for the closed genus-two surface}
The $P_2(2,0)$ action diagram has six vertex orbits, nine edge
orbits and four triangle orbits.  Its barycentric complex of
groups has $19$ named objects, $42$ face homomorphisms and
$24$ coherence triangles.  All local presentations, all $178$
generator images and all twisting elements are given by the
formulas below and the complete certificate
\path{code/d_p2_genus_two.json}.  Exactly one coherence triangle
has a nonidentity twisting element.

The resulting fundamental-group presentation is identified,
on every named generator, with the five-chain presentation (E.27) of
$\Gamma_{2,0}$, verified in E.14.  The inverse sends each of its five
chain generators to a specified local generator.  The unique
dimension-three stratum is the open $\mathcal M_2$; adjoining
it supplies the $24$ explicit cone-coherence identities recorded
in the same certificate.  These are identities of the actual
complex of groups.  A further expansion of these identities
as Peiffer sequences in a compact global presentation is not
asserted here; that additional comparison remains in D.24.
\end{result}
\begin{proof}
Write $\sigma_1,\ldots,\sigma_5$ for the chain twists, and put
\[
\begin{gathered}
 a=\sigma_1,\quad \beta=\sigma_2,\quad b=\sigma_3,
 \quad\gamma=\sigma_4,\quad c=\sigma_5,\\
 W=\sigma_1\sigma_2\sigma_3\sigma_4\sigma_5^2
       \sigma_4\sigma_3\sigma_2\sigma_1,\qquad
 Q=(\sigma_1\sigma_2\sigma_3\sigma_4\sigma_5)^3,\\
 h=(a\beta)^3,\qquad h'=(\gamma c)^3,\qquad
 \delta=h^2=(h')^2,\qquad W=h(h')^{-1}.
\end{gathered}                                                    \tag{D.11}
\]
The global relations make $W$ central of order two and
$Q^2=1$.  For three consecutive chain letters define
$\Delta_4(x,y,z)=xyzxyx$, and set
\[
\begin{aligned}
 t&=\Delta_4(a,\beta,b),&
 e&=a^{-1}cb^{-1}(b\gamma)^3,\\
 u&=tc^{-1},&v&=\Delta_4(b,\gamma,c)a^{-1},\\
 T&=(a\beta)^{-1}(\gamma c)Q.
\end{aligned}                                                       \tag{D.12}
\]
All these are literal words in the same five global generators.
In the machine certificate the local letters $s,d,q$ in the
separating stabilizer mean respectively $\beta,\gamma,Q$;
the local letter $d$ in the nonseparating stabilizer means
$\gamma$.  The separating curve is always denoted $\delta$
in this proof.

Choose the following multicurves:
\[
 N=\{a\},\quad S=\{\delta\},\quad B=\{a,b\},\quad
 L=\{a,\delta\},\quad
 \Theta=\{a,b,c\},\quad D=\{a,c,\delta\}.                    \tag{D.13}
\]
Cutting a genus-two surface shows these are all the orbits:
there are two at each positive rank $1,2,3$.  The independent
stable-graph enumeration gives exactly the counts
$(1,2,2,2)$ at ranks $0,1,2,3$, as checked in the program.

Here are the six stabilizers, with their embeddings fixed
by (D.11)--(D.12).  For $N$ one has
\[
 G_N=\left(
 \bigl(B_4(b,\gamma,c)\times\langle a\rangle\bigr)
 /\langle(b\gamma c)^4=a^2\rangle\right)
 \times C_2(W).                                                \tag{D.14}
\]
To prove this, cutting $a$ gives a torus with two labelled
boundary circles.  Its boundary-framed group has the
four-chain relation $(b\gamma c)^4=z_1z_2$; sewing the two
circles sends $z_1,z_2$ to the same twist $a$.  This is the
orientation-preserving stabilizer of $a$.  The hyperelliptic
involution $W$ reverses the orientation of $a$, as is also
seen from its action $-I$ on first homology.  It is central
of order two and hence supplies exactly the displayed
direct factor, with no unspecified lift.

For $S$, its component-preserving subgroup is
\[
 G_S^+=
 (B_3(a,\beta)\times B_3(\gamma,c))
       /\langle h^2=(h')^2\rangle.
\]
The full group is $G_S=G_S^+\rtimes C_2(Q)$, where $Q$
sends $a\mapsto\gamma$, $\beta\mapsto c$ and exchanges
the two factors.  Sewing the two one-boundary tori proves
the quotient and completeness.  The elements $W$ and
$h(h')^{-1}$ both restrict to $(-I,-I)$ on the two tori;
their ratio is a power of $\delta$.  Both have order two,
so the infinite order of $\delta$ makes that power zero.
This proves their equality in (D.11), including its central
lift.  The braid-shift identity proves the stated action
of $Q$, and the global full-twist relation proves $Q^2=1$.

Put
\[
 H=\bigl(\mathbb Z^2(a,b)\times
          (\mathbb Z(c)*C_2(e))\bigr)\times C_2(W).
\]
This is the subgroup of $G_B$ preserving each of $a,b$.
The full group is obtained by adjoining $t$, with
\[
\begin{gathered}
 tat^{-1}=b,\qquad tbt^{-1}=a,\qquad tct^{-1}=c,
 \qquad tWt^{-1}=W,\\
 tet^{-1}=Wcec^{-1},\qquad t^2=c^2.                            \tag{D.15}
\end{gathered}
\]
For completeness, apply the irreducible stabilizer calculation
of B.9--B.11 inside the four-chain group in (D.14).
Before sewing its branch-switch element
$r=cb^{-1}(b\gamma)^3$ satisfies $r^2=z_1z_2$.
After $z_1=z_2=a$, the element $e=a^{-1}r$ has square one.
This gives precisely $H$, including the independent
orientation-reversing factor $W$ from (D.14).
The element $t$ exchanges $a,b$, so $H$ has index two in
$G_B$.  Relations (D.15), with the specified square and
conjugation action, determine this extension.  Their
possibly nontrivial central lifts were checked below.

The remaining three groups are
\[
\begin{split}
 G_L&=\mathbb Z(a)\times B_3(\gamma,c)\times C_2(W),\\
 G_\Theta&=(\mathbb Z^3(a,b,c)\rtimes\mathfrak S_3(u,v))
                    \times C_2(W),\\
 G_D&=(\mathbb Z^3(a,c,h)\times C_2(W))\rtimes C_2(T).
\end{split}                                                       \tag{D.16}
\]
Here $u$ exchanges $a,b$, $v$ exchanges $b,c$, and
$u^2=v^2=(uv)^3=1$.  The element $T$ exchanges $a,c$
and sends $h\mapsto Wh$, $W\mapsto W$.
The meridian lattice in $G_D$ is $\langle a,c,h^2\rangle$.
For $L$, the left torus stabilizer is
$\langle a,h\rangle$; substituting $h=Wh'$ in the sewing
quotient gives its displayed product.  An element exchanging
the two tori cannot preserve $L$.
For $\Theta$, cutting gives two pairs of pants, with twist
kernel $\mathbb Z^3$ and graph group
$\mathfrak S_3\times C_2$.  The actual words $u,v,W$
supply its displayed splitting.  For $D$, the component-
preserving group comes from
$\langle a,h\rangle\times\langle c,h'\rangle$ with
$h^2=(h')^2$; it is exactly the abelian group in (D.16).
The correcting factors in $T$ are essential: the word $Q$
alone sends $a$ to $\gamma$ and does not preserve (D.13).
The word $T$ has square one and has the indicated action,
so it supplies the full component-exchange extension.

Here is a complete cell-group table.  Define
\[
\begin{gathered}
 A_D=\langle a,c,h,W\rangle=\mathbb Z^3\times C_2,
 \qquad Z_\Theta=\langle a,b,c,W\rangle=\mathbb Z^3\times C_2,\\
 G_{\Theta B}=\langle a,b,c,u,W\rangle,
 \qquad G_{\Theta N}=\langle a,b,c,v,W\rangle.
\end{gathered}
\]
\begin{center}
\begin{tabular}{ll}
\toprule
Cell & Stabilizer\\
\midrule
$N,S,B,L,\Theta,D$ & $G_N,G_S,G_B,G_L,G_\Theta,G_D$\\
$BN$ & $H$\\
$LN,LS$ & $G_L$\\
$\Theta B,\Theta N$ & $G_{\Theta B},G_{\Theta N}$\\
$DB,DS$ & $G_D$\\
$DL,DN$ & $A_D$\\
$\Theta BN$ & $Z_\Theta$\\
$DBN,DLN,DLS$ & $A_D$\\
\bottomrule
\end{tabular}
\end{center}
The curve-permutation actions in (D.15)--(D.16) show both
that these flags exhaust the orbits and that each listed
subgroup is the exact flag stabilizer.  For example,
$\Theta B$ allows the transposition of $a,b$, whereas
$\Theta N$ allows the transposition of $b,c$.  The flag
$\Theta BN$ allows neither.

All chosen flags are nested on (D.13), except that the
$B$ face of $D$ is $\{a,c\}$.  The word $v$ carries this
pair to the chosen $\{a,b\}$ and fixes $a$.  On the
barycentric face category choose transporters
\[
 q_{DB,B}=q_{DBN,BN}=q_{DBN,B}=v,
 \qquad q_{F,F'}=1\text{ on every other face arrow}.          \tag{D.17}
\]
The face homomorphism is, literally,
$x\mapsto q_{F,F'}xq_{F,F'}^{-1}$ in the named target group.
All unchanged named generators map to themselves.  The
nontrivial substitutions needed to specify the other
images are
\[
\begin{array}{ll}
 BN\to N:& e\mapsto a^{-1}cb^{-1}(b\gamma)^3,\\
 \Theta B\to B:&u\mapsto c^{-1}t,\\
 \Theta N\to N:&v\mapsto\Delta_4(b,\gamma,c)a^{-1},\\
 DB\to B:&a\mapsto a,\ c\mapsto b,\
 h\mapsto Wabc^{-1}e,\\
 &T\mapsto ab^{-1}c^{-1}t,\ W\mapsto W,\\
 DL\to L,\ DN\to N:&h\mapsto W(\gamma c)^3,\\
 DS\to S:&h\mapsto(a\beta)^3,\
 T\mapsto(a\beta)^{-1}(\gamma c)Q.
\end{array}                                                       \tag{D.18}
\]
For the higher flag groups use these same substitutions
on their named generators; the full $42$-arrow list in
the JSON certificate specifies every such face individually.
In particular, the factor $ab^{-1}$ in the image of $T$
under $DB\to B$ is retained.

The coherence convention is
\[
 \psi_{F'F''}\psi_{FF'}
       =\operatorname{Ad}(f_{FF'F''})\psi_{FF''},\qquad
 f_{FF'F''}=q_{F',F''}q_{F,F'}q_{F,F''}^{-1}.
\]
There is exactly one nonidentity value:
\[
                  f_{DBN,BN,N}=v\in G_N.                   \tag{D.19}
\]
The fact that $v$ fixes the curve $a$ does not make its
conjugation action on $G_N$ the identity.  Through $BN$
the image of $A_D$ is $vA_Dv^{-1}$; through $DN$ it is
$A_D$.  Equation (D.19) is the exact comparison.

We describe the executed verification and its faithfulness.
The hyperelliptic quotient gives
$\Gamma_2/\langle W\rangle\simeq\Gamma_{0,[6]}$.
Its action on the free group of the six-punctured sphere,
modulo inner automorphisms, is faithful.  The ordinary
Artin action on six letters computes it after eliminating
the last peripheral generator by the product-one relation.
A word with trivial sphere action consequently represents
either $1$ or $W$ in $\Gamma_2$.
The symplectic action distinguishes these two elements,
since $W$ acts by $-I$.
The programs \path{code/b_genus2_checks.py} and
\path{code/b_genus2_faces.py} execute these two exact
checks for all central-lift identities, all $178$ face
images and all $24$ coherence triangles.
The additional program \path{code/d_p2_genus_two.py}
checks all $134$ defining relators in the $11$ local
presentations and records their full words.  Thus no
central $W$ discrepancy is discarded by a sphere-only
calculation.

The face arrows with transporter $1$ contain a connected
spanning graph.  The executed breadth-first construction
records a maximal tree of $18$ such edges.  With this
tree, the global comparison sends every local generator
to its displayed chain word and every face letter to its
transporter (D.17).  The inverse on the five compact
generators is
\[
 \sigma_1\mapsto a_N,\quad
 \sigma_2\mapsto\beta_S,\quad
 \sigma_3\mapsto b_N,\quad
 \sigma_4\mapsto\gamma_N,\quad
 \sigma_5\mapsto c_N.                                       \tag{D.20}
\]
These are literal generator words, with no omitted base
paths.  The action is on the simply connected $P_2(2,0)$
of D.6, and the groups, face maps and twists above are its
actual stabilizer data.  The van Kampen presentation for
this action therefore has the exact comparison
$\pi_1(\text{displayed complex of groups})\to\Gamma_2$
with kernel $\pi_1(P_2)=1$.
It is surjective by (D.20), hence an isomorphism.
Consequently (D.20) respects every compact defining
relation and is inverse on every generator, proving the
claimed full group-presentation dictionary.

Finally adjoin the empty multicurve $O$, representing
the open dimension-three stratum.  Its stabilizer is
$\Gamma_2$ and every transporter into it is $1$.
For each of the $24$ flags $F>F'>F''$ above, the
resulting tetrahedral coherence is exactly
\[
 q_{F',F''}q_{F,F'}
        =f_{FF'F''}\,q_{F,F''}\quad\text{in }\Gamma_2.       \tag{D.21}
\]
The certificate lists both sides as five-chain words;
all $24$ pairs are freely equal words already before
using the global defining relations.  The local groups
and their presentation cells remain part of this action
diagram.  The additional translation of the resulting
higher cellular identities into Peiffer sequences for a
chosen compact presentation is precisely the further
scope retained as open in D.24.
\end{proof}


\subsection{Explicit progress on the three-pointed torus}

\begin{result}{D.28}{Proved}{The exact word problem, flag index and four stabilizer families for $(1,3)$}
Write $a_1,a_2,a_3$ for three mutually disjoint parallel
nonseparating curves, with one marked point in each intervening
annulus, and let $b$ meet each once.  Label $p_i$ as the point
opposite $a_i$ in the standard star picture.  There is the
explicit presentation
\[
 \Gamma_{1,3}=
 \left\langle a_1,a_2,a_3,b\ \middle|\
 [a_i,a_j]=1,\ a_i b a_i=b a_i b,\
 \Delta=(a_1a_2a_3b)^3=1\right\rangle.                       \tag{D.22}
\]
The maps send the four letters to exactly these Dehn twists.
The executed classical Garside algorithm in
\path{code/d_artin_d4.py} solves its word problem by computing
in $A(D_4)$ and retaining the proper greedy factors after
removing the central $\Delta$ exponent.

For $\{i,j,k\}=\{1,2,3\}$, put
\[
 d_i=(a_jba_k)^4,\qquad q_i=a_jba_k a_jba_j.
\]
Then $d_i$ is the separating twist enclosing $p_j,p_k$,
and the following are exact stabilizers with literal
embeddings into (D.22):
\[
\begin{split}
 G_{\{d_i\}}&=B_4(a_j,b,a_k),\\
 G_{\{a_j,a_k\}}&=
 \langle a_1,a_2,a_3,q_i\mid [a_r,a_s]=1,
 q_i a_jq_i^{-1}=a_k,\ q_i a_kq_i^{-1}=a_j,
 (q_i a_i)^2=a_ja_k\rangle,\\
 G_{\{a_j,a_k,d_i\}}&=
 \langle a_j,a_k,q_i\mid[a_j,a_k]=1,
 q_i a_jq_i^{-1}=a_k,\ q_i a_kq_i^{-1}=a_j\rangle,\\
 G_{\{a_1,a_2,a_3\}}&=\mathbb Z^3(a_1,a_2,a_3).
\end{split}                                                       \tag{D.23}
\]
In both middle lines $q_i^2=d_i$; in the third line this
is the separating meridian, not an imposed order-two
relation.  These four families give ten of the twenty-two
vertex-orbit stabilizers.

The complete finite flag index has $22$ vertices, $53$ edges
and $33$ triangles.  Thus the barycentric complex of groups
has $108$ objects, $304$ face arrows and $198$ coherence
triangles.  Every orbit and its actual retained-edge subsets
is listed in \path{code/d_p2_elliptic_three_partial.json}.
D.29 supplies all the remaining named local maps and
completes the full dictionary.  This result records the
exact word algorithm and the four foundational families.
\end{result}
\begin{proof}
Soroko's Proposition~8 and Corollary~9, including their
complete proofs on pages~3--5, identify the boundary-framed
three-holed torus with
\[
 (A(D_4)\times\mathbb Z^3(z_1,z_2,z_3))
       /\langle z_1z_2z_3=\Delta\rangle.
\]
Successive capping kills exactly the three boundary twists
and gives (D.22).  The proof eliminates the other Gervais
generators by the degenerate star relations; in particular
the pair twist opposite $a_i$ becomes $(a_jba_k)^4$.
The figure and those identifications were opened and checked
\cite[Proposition~8 and Corollary~9]{D_Soroko2020}.

Here is the exact algorithm used for subsequent word checks.
Realize the Coxeter group $W(D_4)$ as the $192$ even signed
permutations of four coordinates.  Breadth-first enumeration
supplies the length and a reduced word for each simple element.
The longest element is $-I$, of length $12$, and its positive
lift is the central $\Delta$.  For adjacent simple factors
$x,y$, set
\[
 c=\gcd_{\rm left}(x^{-1}\Delta,y),\qquad
 (x,y)\longmapsto(xc,c^{-1}y).
\]
Each product is length-additive, so this is a positive-monoid
equality.  Repeated transfers terminate: the sequence of
factor lengths strictly increases lexicographically while
its total length is fixed.  Remove identity factors and
central $\Delta$ factors.  The adjacent-gcd criterion is
exactly Dehornoy--Paris Lemma~8.4, whose proof was read;
it identifies the remaining sequence with the unique
greedy normal form \cite[Sections~7--8]{D_DehornoyParis1999}.
A negative atom is replaced by
$\Delta^{-1}(x^{-1}\Delta)$.  Centrality permits collecting
these powers, and cancellativity gives the unique form
$\Delta^m x_1\cdots x_r$ with proper greedy factors.
Two such words agree modulo the central subgroup precisely
when their proper factors agree.

The program computes all $192^2=36864$ left gcds and checks
their uniqueness.  It executes the six Artin relations, four
centrality relations and $400$ seeded contextual rewriting
and normal-form checks.  This supplies a full word algorithm,
not a test in a finite quotient.

Cutting $d_i$ leaves a torus with one marked point and one
boundary, and a pair of pants.  The former has the
four-strand group in (D.23), with boundary twist $d_i$;
this proves the first line, including its meridian.
Cutting $a_j,a_k$ gives a pair of pants and a four-holed
sphere.  The subgroup fixing both curves has kernel
$\mathbb Z^2(a_j,a_k)$ and quotient the pure four-point
group.  The actual twists $a_i,d_i$ are its two standard
intersecting generators.  They commute with the two node
twists and satisfy no further relation, so this subgroup is
\[
             \mathbb Z^2(a_j,a_k)\times F(a_i,d_i).
\]
The only additional graph permutation exchanges the two
node curves.  Its chosen lift is the half twist $q_i$ of
the four-chain subgroup, satisfying $q_i^2=d_i$ and
exchanging $a_j,a_k$.  The star relation gives
\[
                         (q_i a_i)^2=a_ja_k.
\]
The program verifies all these exact words for each of the
three values of $i$; before imposing $\Delta=1$ the last
comparison has precisely one central $\Delta$ factor.
It specifies the conjugation on the remaining free
generator and hence the full index-two extension, proving
the second line of (D.23).

Requiring the additional curve $d_i$ replaces the free
four-point group by its cyclic twist stabilizer.  The
same $q_i$ preserves it, and $q_i^2=d_i$ gives the third
line.  Finally the three complementary pants of
$\{a_1,a_2,a_3\}$ each carry a different fixed label.
Their dual triangle has no nonidentity automorphism
preserving these labels; its stabilizer is exactly its
three independent meridian twists.  This proves the last
line without adding a spurious edge-permutation factor.

For the flag count, use every stable graph of type $(1,3)$
with at least one edge.  For each graph enumerate every
strict chain of nonempty retained-edge subsets, and take
orbits under its full edge-permutation group, realized
by vertex permutations preserving genera and all leg
labels.  Contracting the complementary edges supplies
the actual graph of every face.  This procedure gives
$22,53,33$ orbits in dimensions $0,1,2$.
Each edge has two face arrows; each triangle has six
proper face arrows and six vertex--edge flags.  Hence
there are $2\cdot53+6\cdot33=304$ arrows and
$6\cdot33=198$ coherence triangles.  Parallel-edge
permutations are retained in the enumeration, and
loop half-edge reversals act trivially on these subset
flags.  The saved orbit representatives certify these
counts individually.
\end{proof}


\subsection{The complete three-pointed torus word dictionary}

\begin{result}{D.29}{Proved}{All 108 cells and 304 face maps for $(1,3)$}
The literal $P_2(1,3)$ complex of groups is given by
\path{code/d_p2_elliptic_three.json}; its complete readable
version is \path{code/d_p2_elliptic_three_dictionary.md}.
There are $22$ vertex orbits, $53$ edge orbits and $33$
triangle orbits.  The barycentric diagram has $108$ objects,
$304$ face homomorphisms and $198$ coherence triangles.
Every local group has a named presentation; every generator
image and every twisting element is a word in its named
target generators.  In total the executed certificate checks
$922$ generator images, $126$ defining relators in $28$
local presentations, and $594$ generator equations for the
$198$ coherences.  Exactly seven coherence triangles have
a nonidentity twist.

With the specified $107$-edge maximal tree, its fundamental
presentation is explicitly isomorphic to (D.22).
The inverse sends $a_1,a_2,a_3$ to the corresponding node
generators at the three-curve vertex, and $b$ to the named
$b$ at the all-three-points separating vertex.
The open stratum supplies all $198$ cone tetrahedra,
with the two sides recorded as exact words.  This gives
the fifth requested group-presentation dictionary.
The further compact Peiffer comparisons retained in
D.24 are not asserted by these word identities.
\end{result}
\begin{proof}
We specify the remaining stabilizers, all choices in the
finite diagram, and its inverse presentation map.
Throughout the letters are the actual twists of (D.22).
Put
\[
 h=(a_1b)^3,\quad d_0=h^2,\quad
 \widetilde d_1=d_0d_3^{-1}d_2^{-1},\quad
 \widetilde d_2=d_2,\quad\widetilde d_3=d_3.
                                                               \tag{D.24}
\]
The regular neighborhood of $a_1\cup b$ is a one-holed
torus.  Its complementary four-holed sphere contains
all three marked points.  Soroko's Figure~2 identifies
$d_2,d_3$ with the two standard intersecting pair curves
in that sphere; their intersection number is two.
The boundary-framed lantern relation makes
$\widetilde d_1$ its third pair twist.  The chosen order
in (D.24) is certified by the exact identities
\[
 q_2d_1q_2^{-1}=\widetilde d_1,\qquad
 q_3d_1q_3^{-1}=d_0d_2^{-1}d_3^{-1}.                  \tag{D.25}
\]
Both were checked independently by
\path{code/b_d4_tail.py} using D.28's full word algorithm.
Thus no choice of an unrecorded third pair curve is needed.

Here is the full list of chosen multicurve representatives.
For each $i$, write $\{j,k\}=\{1,2,3\}\setminus\{i\}$
with $j<k$.  In the $L_i$ row use instead the distinguished
ordered pairs
\[
             (j_i,k_i)=(3,2),(1,3),(1,2)
                       \quad(i=1,2,3).
\]
\[
\begin{array}{c|c|c}
\text{name}&\text{multicurve}&\text{number of orbits}\\\hline
N&\{a_1\}&1\\
S_0&\{d_0\}&1\\
S_i&\{d_i\}&3\\
L_0&\{a_1,d_0\}&1\\
L_i&\{a_{j_i},d_i\}&3\\
C_i&\{a_j,a_k\}&3\\
T_i&\{d_0,\widetilde d_i\}&3\\
F_i&\{a_1,d_0,\widetilde d_i\}&3\\
E_i&\{a_j,a_k,d_i\}&3\\
\Theta&\{a_1,a_2,a_3\}&1
\end{array}                                                     \tag{D.26}
\]
Their dual graphs, with every puncture label retained, are
exactly the $22$ nonempty stable graphs of D.28.
The JSON specifies the ordered edge list for every row,
including which edge is the loop and which is the pair
separator.  These orderings determine the masks used below.

\emph{The nonseparating stabilizer.}
Cut $a_1$ and label the four finite points of the resulting
sphere as
\[
         (1,2,3,4)=(\text{node}_1,\text{node}_2,p_3,p_1),
                    \qquad\infty=p_2.
\]
Let $B_4(\mathfrak S_2)$ be the subgroup of $B_4$ whose
permutations belong to $\{1,(12)\}$, with standard pure
generators $A_{ij}$ and the half braid $s=\sigma_1$.
Its presentation is the pure four-braid presentation,
$s^2=A_{12}$, and the six actual conjugation formulas
$sA_{ij}s^{-1}=\sigma_1A_{ij}\sigma_1^{-1}$.
Those right sides are individually expanded in the
certificate, by the exact braid combing of B.18.
Write
$Z_4=A_{12}A_{13}A_{23}A_{14}A_{24}A_{34}$.
Then
\[
 G_N=\bigl(B_4(\mathfrak S_2)\times\mathbb Z(a_1)\bigr)
                   /\langle Z_4a_1^2=1\rangle.          \tag{D.27}
\]
Its literal generator embedding is
\[
\begin{aligned}
 s&=h a_1^{-1},& A_{12}&=s^2,\\
 A_{13}&=a_2a_1^{-1},& A_{14}&=a_2^{-1}a_3d_2^{-1},\\
 A_{34}&=d_2,& A_{23}&=sA_{13}s^{-1},\qquad
 A_{24}=sA_{14}s^{-1}.
\end{aligned}                                                     \tag{D.28}
\]
To prove this identification, use zero-rotation
parameterizations of the two node holes.  A pure braid
is lifted by rigid translation of disjoint sufficiently
small discs; their radii can be chosen continuously from
the minimum point separation.  This gives the specified
zero framing.  Sewing identifies both node-hole boundary
twists with $a_1$.  Capping the outer boundary gives
$Z_4(z_1z_2)=1$, hence precisely (D.27).
The half exchange has image $h a_1^{-1}$; its square is
$d_0a_1^{-2}$.  A twist enclosing node $1$ and $p_3$ has
image $a_2a_1^{-1}$, and the pair $p_3,p_1$ has image
$d_2$.  The three-point cluster $\{1,3,4\}$ has full
twist $A_{13}A_{14}A_{34}$ and image $a_3a_1^{-1}$;
solving yields the stated $A_{14}$.  Conjugation by the
actual half exchange gives the other two generators.
These generate the entire cut-sphere group, and the
only sewing kernel is the displayed boundary relation.
This proves both surjectivity and injectivity of (D.27),
including the branch-exchanging coset.
Useful inverse words, also checked exactly, are
\[
\begin{aligned}
 a_2&=A_{13}a_1,&a_3&=a_2A_{14}A_{34},&h&=sa_1,\\
 d_3&=a_1^2A_{12}A_{13}A_{23},&
 \widetilde d_1&=a_1^2A_{12}A_{14}A_{24}.&&
\end{aligned}                                                     \tag{D.29}
\]

\emph{The remaining vertex groups.}
The four families $S_i,C_i,E_i,\Theta$ are proved in
D.28.  The other groups have the following full
presentations, with literal generator names:
\[
\begin{aligned}
 G_{S_0}&=B_3(a_1,b)\times F(d_2,d_3),\\
 G_{L_0}&=\mathbb Z^2(a_1,h)\times F(d_2,d_3),\\
 G_{T_i}&=B_3(a_1,b)\times\mathbb Z(\widetilde d_i),\\
 G_{F_i}&=\mathbb Z^3(a_1,h,\widetilde d_i),\\
 G_{L_i}&=\mathbb Z(a_{j_i})\times
       \langle a_{k_i},r_i\mid[r_i^2,a_{k_i}]=1\rangle,\\
 r_i&=a_{k_i}a_{j_i}^{-1}(a_{j_i}b)^3,
                       \qquad r_i^2=d_i.
\end{aligned}                                                     \tag{D.30}
\]
For $S_0$, cut the one-holed torus from the three-puncture
disc.  The latter's pure braid group has central boundary
$\mathbb Z(d_0)$ and free complementary factor
$F(d_2,d_3)$.  Sewing identifies $d_0$ with $h^2$ in
$B_3$ and eliminates this one central generator; it gives
exactly the first product.  Stabilizing $a_1$ in its
$B_3$ factor gives $\mathbb Z^2(a_1,h)$, as in B.10.
Stabilizing a specified pair curve in the four-holed
sphere replaces its free group by that curve's cyclic
stabilizer.  This proves the next two products and the
$F_i$ row, with the node meridian $h^2$ retained.
Finally $L_i$ is the boundary-retaining nonseparating
stabilizer inside $B_4(a_{j_i},b,a_{k_i})$.
B.10's cutting proof gives
$h_i=a_{j_i}a_{k_i}^{-1}r_i$ and $r_i^2=d_i$;
it gives the last presentation without imposing a
spurious order-two relation on $r_i$.
The program checks each relation and each displayed
embedding in the exact group (D.22).

\emph{Flag stabilizers and all transports.}
The only nontrivial permutations of distinct edges of
these representatives occur at $C_i,E_i$: $q_i$ exchanges
the two nonseparating edges.  Their ordered subgroups are
\[
 G_{C_i}^{\mathrm{ord}}
    =\mathbb Z^2(a_j,a_k)\times F(a_i,d_i),\qquad
 G_{E_i}^{\mathrm{ord}}=\mathbb Z^3(a_j,a_k,d_i).          \tag{D.31}
\]
A flag has the full group of its largest multicurve
unless one of its subsets distinguishes these two edges;
in that case it has the appropriate ordered subgroup.
Loop-branch reversals in $N,L_i,L_0,F_i$ remain in their
already displayed groups; they do not permute distinct
edges.  Thus (D.27), (D.30), D.28 and (D.31) give all
$28$ presentations used by the $108$ objects.

Here is the completely deterministic representative
rule used in the supplied tables.  For each largest
multicurve, take every decreasing chain of nonempty
edge masks and choose the lexicographically least one
in its actual edge-permutation orbit.  This gives the
$22,53,33$ orbits.  For a face, first take its largest
retained mask.  Among
\[
             1,q_1,q_2,q_3,q_1^{-1},q_2^{-1},q_3^{-1}
\]
choose the first word carrying its individual twist
set to the chosen representative in (D.26).
Such a word exists for every subset in the executed
list.  Transport the remaining masks by its actual
edge bijection.  If they are not the chosen least
flag, multiply on the left by the target $q_i$ that
puts them there.  This is possible precisely in the
$C_i,E_i$ cases.  The resulting word is $u_a$ for the
face arrow $a$.

These are actual transporters.  The identities are
checked by conjugating each twist and using D.28's
full word problem.  Equality of two positive Dehn
twists identifies their curves: if $T_c=T_d$, applying
them to $c$ and using
$i(c,T_d(c))=i(c,d)^2$ first makes the curves disjoint;
independence of the twists of a multicurve then makes
them equal.  The ordered masks consequently identify
the entire face flag, not only its abstract dual graph.
The face injection is exactly
$\psi_a(g)=u_a g u_a^{-1}$.

Every such injection is rewritten in the target's
named generators.  Equations (D.28)--(D.30) provide the
needed inverse substitutions at $N$ and $L_i$.
Conjugations of the remaining leaves are obtained from
\[
 q_i a_i q_i^{-1}=a_ja_k a_i^{-1}d_i^{-1},\qquad
 q_i^{-1}a_iq_i=a_ja_k d_i^{-1}a_i^{-1},
\]
while $q_i$ exchanges $a_j,a_k$ and fixes $b$.
Together with (D.25), these formulas cover every
transport in the list.  Each of the $922$ resulting
words is independently evaluated in (D.22) and
checked against $u_a g u_a^{-1}$.  The faithful
geometric embeddings of the local groups then prove
that these are homomorphisms into the named target
presentations, with no kernel introduced by rewriting.

For composable face arrows $a,b$, the twisting word is
\[
                    f_{b,a}=u_bu_a u_{ba}^{-1},\qquad
             \psi_b\psi_a=\operatorname{Ad}(f_{b,a})\psi_{ba}.
\]
For clarity the nonidentity
values can be listed completely.  A row specifies the
triangle flag and the intermediate edge flag; the
last column gives its target-group element.  The
subsets below are the actual curves in
the largest representative before its recorded face
transport.
\[
\begin{array}{c|c|c|c}
\text{triangle}&\text{intermediate edge}&\text{target}&f_{b,a}\\\hline
(F_1,\{a_1,\widetilde d_1\},\widetilde d_1)
 &(L_1,d_1)&S_1&a_3^{-1}b^{-1}a_3^{-1}a_2ba_3\\
(E_1,\{a_2,d_1\},a_2)&(L_1,a_3)&N&h d_3^{-1}\\
(E_1,\{a_2,d_1\},d_1)&(L_1,d_1)&S_1&q_1\\
(\Theta,C_1,a_3)&(C_1,a_2)&C_1&q_1\\
(\Theta,C_1,a_3)&(C_1,a_2)&N&h d_2^{-1}\\
(\Theta,C_2,a_3)&(C_2,a_1)&C_2&q_2\\
(\Theta,C_3,a_2)&(C_3,a_1)&C_3&q_3
\end{array}                                                     \tag{D.32}
\]
The fourth and fifth rows use the two possible final
vertices of the same intermediate edge.  The words
at $N$ are read by (D.29), and $q_1$ at $S_1$ is its
ordinary four-braid half twist.  Their exact global
checks are
\[
\begin{aligned}
 q_3q_1q_2^{-1}&=h d_2^{-1},&
 q_2q_1q_3^{-1}&=h d_3^{-1},\\
 q_2^{-1}q_3^{-1}&=a_3^{-1}b^{-1}a_3^{-1}a_2ba_3.&&
\end{aligned}
\]
All other $191$ twisting elements are identity.
The $198$ equations are checked on all $594$ source
local generators in the certificate.  This establishes
the entire complex of groups with its actual twists.

\emph{The presentation maps and dimension-three words.}
Select the first spanning tree in the given arrow order
using only $u_a=1$ arrows; the executed list has $107$
edges and connects all $108$ objects.  The resulting
finite presentation has the object-group generators
and arrow letters $t_a$, relations
\[
 t_a g t_a^{-1}=\psi_a(g),\qquad
 t_b t_a=f_{b,a}t_{ba},\qquad
 t_a=1\ \text{on the chosen tree}.
\]
The map sends $t_a$ to $u_a$; together with the literal
local embeddings this defines its map to (D.22).  The diagram is the
actual action complex of groups on $P_2(1,3)$: its
cell representatives and stabilizers were identified
above, including every parallel-edge incidence.
D.6 proves that this development is simply connected.
Lifting edge paths and elementary face homotopies
therefore proves that the displayed fundamental-group
map is an isomorphism.  Its inverse on the four
presented global generators is the one stated in the
result, since their local images are literally
$a_1,a_2,a_3,b$.  The isomorphism proves that these
specified inverse words satisfy all six Artin
relations and the central relation of (D.22).

Finally adjoining the empty multicurve is coning the
whole development.  Set the transporter from every
object to this open vertex equal to $1$.  The twist
on an object--face--open triangle is then precisely
its already specified $u_a$.  For each of the $198$
three-step flags the tetrahedral identity becomes
\[
                       u_bu_a=f_{b,a}u_{ba}.
\]
Both sides agree as freely reduced words, not only as
elements of the presented group.  These are all the
cone tetrahedra for the unique dimension-three
stratum of $M_{1,3}$.  Their further expansion in a
chosen compact presentation's free crossed module
is the additional task preserved in D.24.
\end{proof}


\subsection{An explicit cellular lift of the six-point comparison}

\begin{result}{D.30}{Proved}{The 630 geometric tetrahedra and the 24 compact prism cells}
There is a specified CW model of the full six-point
complex of groups, with its compact open-stratum model
$K_6$ from D.25, and a cellular homotopy equivalence
\[
                       F:\mathcal K\longrightarrow K_6
\]
whose restriction to that open-stratum subcomplex is
literally the identity.  Its two-skeleton uses the exact
local generators, sewing words and transports of D.22.
The complete cellular comparison through those cells is
\path{code/d_p5c3_geometric_comparison.json}.

The source has $236$ vertices, $2194$ oriented one-cells
and $8371$ two-cell relators.  It is a groupoid
presentation; its $235$ crossing edges form a tree.
Collapsing that tree leaves $1959$ one-vertex generators.  Every relator image is supplied as an
explicit product of conjugates of the fixed $26$ compact
relators.  Each of the $630$ geometric cone tetrahedra
has a fully expanded identity sequence in the source
alphabet.  Its image cancels by adjacent inverse
Peiffer factors.  Thus its coordinates in D.25's
$24$-element basis are all zero.  The $24$ open-stratum
prism cells map identically to that basis.  On these
$654$ specified three-cells the exact comparison matrix is
\[
                     \bigl[\,0_{24\times630}\mid I_{24}\,\bigr].
                                                               \tag{D.33}
\]
The inverse on the compact basis is inclusion of the
named open prism cells.  In the original P5c3 orientation,
D.25 changes the second block to $-I_{24}$ and reverses
those inverse prism orientations.

These coordinates refer to the specified cellular
model and its recorded coherence fillings.  The theorem
supplies those choices and their simultaneous
compatibility; it does not assign a canonical
Peiffer class to a bare equality of group words.
\end{result}
\begin{proof}
\emph{A relator-valued normalizer.}
Let $F_9$ be the free group on the nine compact letters,
and let $R_0,\ldots,R_{25}$ be exactly D.25's relators.
For a free word $w$, the program
\path{code/d_p5c3_fill.py} returns its unique semidirect
normal form $\nu(w)$ and a sequence $N(w)$ such that
\[
     \prod_{(c,r,\epsilon)\in N(w)}cR_r^\epsilon c^{-1}
                            =w\nu(w)^{-1}
                   \quad\text{freely in }F_9.             \tag{D.34}
\]
Here $\nu(w)$ is a word in $a,b$, followed by a word in
$x,y,z$, followed by a word in $p,q,r,s$.
The two free-group actions are the actual actions of
B.19, including their inverse automorphisms.

For completeness, suppose the current normal prefix
is $uv$, with $v$ in the final free factor, and append
a letter $h$ from the base.  Put $v'=\phi_{h^{-1}}(v)$.
The change from $uvh$ to $uhv'$ is certified by
\[
       u\,D_h(v')^{-1}u^{-1},\qquad
       D_h(v')=hv'h^{-1}\phi_h(v')^{-1}.
\]
The four expansion rules (D.9) express this as conjugate
mixed relators.  Normalize $uh$ recursively in the
five-generator base, using its six relators.  Appending
a fibre letter needs only free reduction.  Multiplying
these certificates telescopes to (D.34).
Uniqueness follows from the two proved free-by-free
semidirect decompositions.  In particular $w=1$ in
$\Gamma_{0,6}$ exactly when $\nu(w)$ is empty, and then
$N(w)$ is a complete van Kampen certificate for $w$.
All $26$ defining relators, all $52$ signed action-inverse
checks, and $20$ finite contextual checks were executed.
The subsequent construction independently checks (D.34)
for every relator it actually uses.

\emph{The complete source two-skeleton.}
Adjoin the open object $0$ to D.22's diagram.
Use the compact presentation at $0$, a separate named
copy of every local presentation, an arrow letter
$e_a$ for every boundary arrow $a:i\to j$, and a
letter $t_i$ for every crossing $i\to0$.
Write $q_a$ for D.22's compact word of its transporter,
$\mu_i$ for the compact word map on its local generators,
and $f_{b,a}$ for its specified local twisting word.
For a composable pair $a:i\to j$, $b:j\to k$, put
\[
\begin{aligned}
 S_{b,a}&=e_b e_a e_{ba}^{-1}f_{b,a}^{-1},\\
 O_a&=t_j e_a t_i^{-1}q_a^{-1},\\
 C_i(x)&=t_i x t_i^{-1}\mu_i(x)^{-1}.
\end{aligned}                                                     \tag{D.35}
\]
Include the local relators, the ordinary sewing
relators, all $C_i(x)$ for named local generators,
all $S_{b,a}$, and all $O_a$.
The counts, each obtained from the actual lists, are
\[
\begin{array}{c|r}
\text{relator family}&\text{number}\\\hline
\text{open compact}&26\\
\text{boundary local}&2195\\
\text{boundary sewing}&3570\\
\text{local-generator crossing }C_i(x)&1110\\
\text{boundary triangle }S_{b,a}&630\\
\text{open triangle }O_a&840
\end{array}
\]
Their sum is $8371$; the $9+1110+840+235=2194$
letters are the oriented one-cells just specified.
The $235$ boundary objects and the open object are
its vertices.  The $t_i$ are a spanning tree, not
independent generators after passing to a one-vertex
fundamental-group presentation.
No local relator or sewing map is omitted from this
source two-presentation.

Define its map on free letters by
\[
  x_i\longmapsto\mu_i(x_i),\qquad
  e_a\longmapsto q_a,\qquad t_i\longmapsto1,
\]
and use the identity on the nine open letters.
For each relator, use (D.34) as its two-cell filling;
for an open compact relator use its own single
relator cell.  These choices agree because the
normalizer returns that same certificate there.
The crossing relators $C_i(x)$ and open triangles
$O_a$ have freely empty images and get empty fillings.
The full list uses $1027$ distinct certificates with
$26483$ conjugate-relator factors.  Every expanded
product was checked against the image of its source
relator in the free group $F_9$.
These free-word certificates prove the chain-boundary
identity in degree two, including all local and
sewing columns: applying the Fox boundary to either
side of (D.34) gives exactly the same lifted edge chain.

\emph{The actual tetrahedral identity sequences.}
For a local word $z$, expand $C_k(z)$ as a sequence
in the generator crossing relators $C_k(x)$.
For a positive letter its conjugator is the image
under $\mu_k$ of the preceding prefix; for a negative
letter multiply that conjugator by $\mu_k(x)^{-1}$
and use $C_k(x)^{-1}$.  Free cancellation proves that
this sequence has boundary
$t_kzt_k^{-1}\mu_k(z)^{-1}$.
Thus this notation specifies a finite relator sequence,
not an extra generator or an unspecified equality.

Fix one of the $630$ flags $i\to j\to k\to0$, and
write $a:i\to j$, $b:j\to k$, $f=f_{b,a}$,
$\mu f=\mu_k(f)$.  Let $T_{b,a}$ be the already
chosen compact certificate for the image of
$S_{b,a}$.  Its boundary is
\[
             w_{b,a}=q_bq_aq_{ba}^{-1}(\mu f)^{-1}.
\]
The actual local coherence at the open vertex is
\[
             \delta_{b,a}=\mu f\,q_{ba}(q_bq_a)^{-1}
                          =w_{b,a}^{-1}.
\]
Choose its open-stratum van Kampen diagram to be
$T_{b,a}^{-1}$, with the same $26$ relator cells and
reversed orientations.  This is a single specified
rule for all flags.  The complete tetrahedral
identity sequence in the full source alphabet is
\[
\begin{split}
 &(t_k S_{b,a}t_k^{-1})\ C_k(f)\
       (\mu f\,O_{ba}(\mu f)^{-1})\ T_{b,a}^{-1}\\
 &\hspace{25mm}\cdot
       (q_bO_aq_b^{-1})^{-1}\ O_b^{-1}=1.
\end{split}\tag{D.36}
\]
The notation $T_{b,a}^{-1}$ in this source expression
means the actual sequence of open compact relators;
their indices and conjugators are recorded individually.

Here is a direct verification of the attaching word.
Set $W=t_ke_be_at_i^{-1}$.
The route through the two open triangles gives
\[
 W=O_b\,(q_bO_aq_b^{-1})\,q_bq_a.
\]
The route through the boundary triangle gives
\[
 W=(t_kS_{b,a}t_k^{-1})\ C_k(f)\
             (\mu f\,O_{ba}(\mu f)^{-1})\,\mu f\,q_{ba}.
\]
Comparing them and inserting the specified diagram
for $\delta_{b,a}$ yields (D.36) by free cancellation.
This is precisely the tetrahedral comparison of the
two boundary routes, including its local open coherence.
The program expands every factor back into the full
$2194$-letter source alphabet and obtains the empty
word for each of the $630$ sequences.

Apply the two-map factor by factor to (D.36).
The first factor gives $T_{b,a}$ since $t_k$ maps to
$1$.  Each crossing and open-face factor has an
empty certificate.  The open coherence gives its
literal inverse $T_{b,a}^{-1}$.  The resulting
identity sequence is therefore
\[
                         T_{b,a}T_{b,a}^{-1}.             \tag{D.37}
\]
Cancel its adjacent inverse conjugate-relator factors
successively.  These are explicit Peiffer cancellations,
so its class in $\pi_2(K_6^{(2)})$ is zero, not merely
its boundary word or its augmentation.
Only two of the $630$ certificates $T_{b,a}$ are
nonempty; the eight mapped factors across all
$630$ sequences were checked individually.
The source sequences contain $2539$ factors in all.
This proves the zero block of (D.33) at the level of
actual identity sequences.

\emph{Simultaneous CW realization and the inverse.}
There is an explicit mapping-cylinder construction,
which fixes the choices for all cells at once.
Let $X=|P_2(0,6)|$ and
$L=E\Gamma\times_\Gamma X$, where
$\Gamma=\Gamma_{0,6}$.  Give $L$ the classifying
complex-of-groups CW model of D.22.  It has $235$
vertices, $1950$ one-cells and $6395$ two-cells:
the latter are its $2195$ local, $3570$ sewing and
$630$ triangle cells.  Higher cells include the
local classifying-space cells and their prisms.
Its projection $p:L\to B\Gamma$ is the actual
classifying map of its universal development.
There is a direct identification
\[
 E\Gamma\times_\Gamma\operatorname{Cone}(X)
                   \cong M_p.
\]
Indeed $[e,(x,t)]$ corresponds to $([e,x],t)$
away from the cone vertex, and its endpoint is
$[e]\in B\Gamma$.  Thus the full open-stratum
cone diagram is the mapping cylinder of this
specific projection.

Replace $B\Gamma$ by D.25's $K_6$ and choose the
cellular map $f:L\to K_6$ whose one- and two-cell
maps are exactly the words and certificates already
specified.  It extends to all higher cells because
$K_6$ is aspherical.  Its induced fundamental-group
map is the one of $p$, so it is homotopic to $p$
under $B\Gamma\simeq K_6$: lift the comparison to
their contractible universal targets and extend it
cell by cell.  Set $\mathcal K=M_f$.
This is a CW replacement of the actual cone
classifying diagram, with its boundary and open
maps specified.  The map
\[
 F|_{K_6}=1,\qquad F(z,t)=f(z)\quad(z\in L)
                                                               \tag{D.38}
\]
is its canonical cellular retraction.

Its two-skeleton is precisely (D.35): the vertical
edges are the $t_i$, the prisms of local generator
edges are $C_i(x)$, and the prisms of boundary
arrows are $O_a$.  A prism over a boundary two-cell
uses its one fixed compact filling as its top
face.  In particular the prism over $S_{b,a}$
uses $T_{b,a}^{-1}$ with the opposite orientation.
Its side annulus expands into the crossing and
open-face terms of (D.36).  This proves that the
$630$ coherence disks are the restrictions of
one mapping-cylinder construction.  They have
not been varied independently after fixing their
boundaries.

The program also constructs the standard prism
identity for every one of the $6395$ boundary
two-cells.  For a cell relator $R$ based at $i$,
let $C(R)$ be the word obtained by expanding its
side annulus in the local-generator and arrow
crossing relators.  Its free boundary is
$t_iRt_i^{-1}\mu(R)^{-1}$, with the same
prefix rule used above.  The prism identity is
\[
          (t_iRt_i^{-1})\,N(\mu(R))^{-1}\,C(R)^{-1}=1.
                                                               \tag{D.39}
\]
All $6395$ source identities and their factorwise
images were freely checked.  Each image is its
own compact certificate followed by its inverse.
For a triangle, grouping the side annulus by its
three faces gives (D.36); moving the compact
certificate across that annulus is the defining
Peiffer interchange.  Thus the two expansions
are the same oriented prism attaching sphere.

Equation (D.38) sends every such three-dimensional
prism into $K_6^{(2)}$, proving its degree-three
chain image is zero.  The full source also has
three-cells of $L$, including local three-cells,
edge times local two-cells, and triangle times
local one-cells.  These are not being omitted
from the model or identified with the selected
$654$ summands in (D.33).  Their images are those
of the already extended map $f$.  At chain level
they are uniquely determined by their boundaries,
since D.25 proves that $\partial_3$ is injective.
The computed zero prism block and the identity
open block therefore belong to an actual full
chain map; (D.33) specifies those summands only.

The mapping cylinder retracts onto $K_6$, so
(D.38) is the asserted homotopy equivalence.  Inclusion
$K_6\hookrightarrow\mathcal K$ is a strict right
inverse for the specified map on this subcomplex.
Its $24$ prism attaching sequences are exactly
D.25's sequences.  The program maps each factor
of all $24$ sequences identically and verifies
their freely empty boundaries again.  Hence the
identity block and its inverse in (D.33) are
literal.  The final sign comparison with P5c3
is the already proved individual inversion in
D.25, completing the stated cellular comparison.
\end{proof}

